\documentclass[11pt,oneside]{article}

\usepackage[T1]{fontenc}
\usepackage{lmodern}
\usepackage{microtype}
\usepackage[margin=0.92in,headheight=15pt]{geometry}
\usepackage{amsmath,amssymb,amsthm,mathtools}
\usepackage{booktabs,longtable,tabularx,array,multirow}
\usepackage{enumitem}
\usepackage[dvipsnames]{xcolor}
\usepackage{hyperref}
\usepackage[nameinlink,noabbrev]{cleveref}
\usepackage{fancyhdr}
\usepackage{titlesec}
\usepackage{needspace}
\usepackage{listings}
\usepackage{xurl}
\usepackage{seqsplit}

\definecolor{LeanBlue}{RGB}{28,86,140}
\definecolor{PaperGreen}{RGB}{20,112,72}
\definecolor{NewViolet}{RGB}{99,56,150}
\definecolor{ComputeOrange}{RGB}{166,86,20}
\definecolor{CautionRed}{RGB}{150,42,42}
\definecolor{MutedGray}{RGB}{88,96,105}
\definecolor{CodeBg}{RGB}{247,248,250}

\hypersetup{
  colorlinks=true,
  linkcolor=LeanBlue,
  citecolor=PaperGreen,
  urlcolor=LeanBlue,
  pdftitle={Toward the Alaoglu--Erdos Integer-Exponent Conjecture},
  pdfauthor={OpenAI GPT-5.6 Pro; prepared for Vladimir Reshetnikov},
  pdfsubject={Audit of the ProveIt corpus, new partial results, equivalent formulations, and research directions}
}
\urlstyle{same}

\pagestyle{fancy}
\fancyhf{}
\lhead{Alaoglu--Erd\H{o}s integer-exponent conjecture}
\rhead{Research report}
\cfoot{\thepage}
\renewcommand{\headrulewidth}{0.4pt}

\titleformat{\section}{\Large\bfseries}{\thesection}{0.72em}{}
\titleformat{\subsection}{\large\bfseries}{\thesubsection}{0.72em}{}
\titleformat{\subsubsection}{\normalsize\bfseries}{\thesubsubsection}{0.72em}{}

\setlist[itemize]{leftmargin=2em,itemsep=0.26em,topsep=0.35em}
\setlist[enumerate]{leftmargin=2.25em,itemsep=0.26em,topsep=0.35em}
\setlength{\parindent}{0pt}
\setlength{\parskip}{0.52em}
\setlength{\emergencystretch}{3em}
\allowdisplaybreaks

\newtheorem{theorem}{Theorem}[section]
\newtheorem{proposition}[theorem]{Proposition}
\newtheorem{lemma}[theorem]{Lemma}
\newtheorem{corollary}[theorem]{Corollary}
\newtheorem{conjecture}[theorem]{Conjecture}
\newtheorem{question}[theorem]{Question}
\theoremstyle{definition}
\newtheorem{definition}[theorem]{Definition}
\newtheorem{observation}[theorem]{Observation}
\newtheorem{problem}[theorem]{Research target}
\theoremstyle{remark}
\newtheorem{remark}[theorem]{Remark}
\newtheorem{warning}[theorem]{Caution}

\newcommand{\N}{\mathbb N}
\newcommand{\Nzero}{\mathbb N_0}
\newcommand{\Z}{\mathbb Z}
\newcommand{\Q}{\mathbb Q}
\newcommand{\R}{\mathbb R}
\newcommand{\C}{\mathbb C}
\newcommand{\Qbar}{\overline{\mathbb Q}}
\newcommand{\Apos}{\overline{\mathbb Q}_{>0}}
\newcommand{\Sol}{\mathcal S}
\newcommand{\Cand}{\mathcal C}
\newcommand{\GammaSat}{\Gamma_{\mathrm{sat}}}
\newcommand{\LL}{\widetilde{\mathcal L}}
\newcommand{\vpp}[2]{v_{#1}\!\left(#2\right)}
\newcommand{\distZ}[1]{\left\lVert #1\right\rVert_{\Z}}
\newcommand{\Span}{\operatorname{span}}
\newcommand{\rank}{\operatorname{rank}}
\newcommand{\supp}{\operatorname{supp}}
\newcommand{\PGL}{\operatorname{PGL}}
\newcommand{\sha}[1]{\texttt{\seqsplit{#1}}}
\newcommand{\lean}[1]{\path{#1}}
\newcommand{\statuslean}{\textcolor{LeanBlue}{\textbf{[Lean]}}}
\newcommand{\statuspaper}{\textcolor{PaperGreen}{\textbf{[corpus paper proof]}}}
\newcommand{\statusnew}{\textcolor{NewViolet}{\textbf{[new paper proof]}}}
\newcommand{\statuscomp}{\textcolor{ComputeOrange}{\textbf{[interval computation]}}}
\newcommand{\statusclassical}{\textcolor{MutedGray}{\textbf{[classical]}}}
\newcommand{\statusopen}{\textcolor{CautionRed}{\textbf{[open]}}}

\newenvironment{statusbox}[1]
{\begin{center}\begin{minipage}{0.95\textwidth}\raggedright\hrule\vspace{0.45em}\textbf{Status.} #1\par\vspace{0.35em}}
{\vspace{0.35em}\hrule\end{minipage}\end{center}}

\lstdefinestyle{leanstyle}{
  basicstyle=\ttfamily\footnotesize,
  backgroundcolor=\color{CodeBg},
  frame=single,
  rulecolor=\color{MutedGray},
  showstringspaces=false,
  breaklines=true,
  columns=fullflexible,
  keepspaces=true,
  tabsize=2
}

\begin{document}
\pagenumbering{gobble}

\begin{titlepage}
\centering
\vspace*{0.9cm}
{\Huge\bfseries Toward the Alaoglu--Erd\H{o}s\par}
\vspace{0.2cm}
{\Huge\bfseries Integer-Exponent Conjecture\par}
\vspace{0.65cm}
{\LARGE Corpus Audit, Strong Algebraic-Base Spectra,\par}
\vspace{0.12cm}
{\LARGE Tower Rigidity, and a Branch-Sensitive Research Program\par}
\vspace{0.95cm}
{\large A detailed analysis of\par}
\vspace{0.1cm}
{\large \url{https://github.com/VladimirReshetnikov/ProveIt}\par}
\vspace{0.2cm}
{\normalsize at repository commit\par}
\vspace{0.08cm}
{\normalsize\sha{4848c6d834190301e014028a9c78a3f0ef49b7b8}\par}
\vfill
\begin{minipage}{0.91\textwidth}
\small
\textbf{Verdict.}
I did not obtain an unconditional proof of the conjecture.  The obstruction is not a
missing elementary lemma: a hypothetical counterexample supplies precisely the unresolved
$2\times2$ logarithmic determinant underlying the four exponentials conjecture.  The
repository has nevertheless moved well beyond that bare reduction.  Its Lean development
formalizes a sharp conditional semigroup structure, finite exclusions, arithmetic rigidity,
counting and equidistribution statements, radical reformulations, a generic algebraic
six-exponentials determinant, and the exact positive-algebraic output locus.

At the audited snapshot, Lean already proves that under a nonintegral solution the positive
algebraic bases with algebraic $x$-th power are exactly
$\Gamma=\{2^r3^s:r,s\in\Q\}$.  This report gives a short general two-base derivation of
that phenomenon and adds results not isolated in the current article corpus: Roy's strong
six-exponentials upgrade; a Smith-normal-form lattice for rational and integral outputs; a
solution-independent second-iterate kernel; its rational M\"obius-orbit characterization;
depth-three power-tower rigidity; and new one-base tower equivalences.  Except where marked
\statuslean, these additions are complete paper proofs, not yet machine-checked or
peer-reviewed.
\end{minipage}
\vfill
{\large Prepared for Vladimir Reshetnikov\par}
\vspace{0.12cm}
{\normalsize OpenAI GPT-5.6 Pro\par}
\vspace{0.35cm}
{\large August 20, 2026\par}
\end{titlepage}

\pagenumbering{roman}
\begin{abstract}
The Alaoglu--Erd\H{o}s integer-exponent conjecture, originating in the 1944 work of Alaoglu and Erd\H{o}s~\cite{AlaogluErdos1944}, asserts that a real number $x$ for which
$2^x$ and $3^x$ are integers must itself be an integer.  I audited all three LaTeX reports
under \texttt{Analysis/ExponentialIdentities/Article/**/} and the current Lean import surface
at the commit printed on the title page.  The corpus proves, among other things, that a
counterexample is transcendental; that failure produces a unique least nonintegral solution
$\beta$ and the exact solution monoid $\Nzero+\Nzero\beta$; that the corresponding candidate
integers form a free rank-two monoid; that nonintegral candidates have zero density and obey
strong progression and determinant repulsion estimates; and that the algebraic-output locus
among all positive algebraic auxiliary bases is exactly the positive $2,3$-unit divisible
hull, with additional mixed-second-iterate rank restrictions.  The Lean kernel excludes
nonintegral solutions with $0<x<12$, while a separate directed-rounding MPFR computation
reported in the third article excludes nonintegral solutions with $0<x\le 26$; the latter is
source-audited here but not reclassified as a formal proof.

A central theorem of the current Lean snapshot is the exact $2,3$ algebraic-base spectrum:
under a nonintegral solution, a positive algebraic $\alpha$ has algebraic $\alpha^x$ exactly
when some positive power of $\alpha$ is a rational $2,3$-unit, equivalently when
$\alpha\in\Gamma=\{2^r3^s:r,s\in\Q\}$.  I give a short general two-base version.  If
$a,b$ are multiplicatively independent positive algebraic numbers, $x$ is irrational, and
$a^x,b^x$ are algebraic, then
\[
  \alpha^x\in\Qbar
  \quad\Longleftrightarrow\quad
  \alpha\in\GammaSat(a,b):=\{a^r b^s:r,s\in\Q\}.
\]
The difficult direction is an exact application of six exponentials and is compatible with
the repository's generic algebraic determinant.  The genuinely new upgrade uses Baker's
homogeneous theorem and Roy's strong six exponentials theorem to prove
$x\log\alpha\notin\LL$ outside the saturated group, where $\LL$ is the $\Qbar$-linear
span of $1$ and logarithms of algebraic numbers.  Consequently every fixed algebraic
$\alpha\notin\Gamma$ is a complete one-number detector for the conjecture, generalizing
the repository's base-$5$ equivalence.

For a hypothetical nonintegral solution, put $M=2^x$ and $A=3^x$.  Inside $\Gamma$, rational
and integral outputs admit the exact valuation criteria
\[
  (2^r3^s)^x=M^rA^s\in\Q
  \Longleftrightarrow
  r\,\nu(M)+s\,\nu(A)\in\Z^{(\mathbb P)},
\]
with nonnegativity added for integrality.  This gives a finite-index rank-two superlattice of
$\Z^2$, whose quotient is read off from the Smith invariants of the prime-valuation matrix.
A second, coarser kernel
\[
 K_x=\{(r,s)\in\Q^2:M^rA^s\in\Gamma\}
\]
is independent of the chosen nonintegral solution.  It has dimension at most one, contains
no vector with $rs>0$, and is nonzero exactly when
\[
  x=\frac{u+v\vartheta}{r+s\vartheta},
  \qquad \vartheta=\log_2 3,
\]
for rational coefficients with nonzero determinant.  Thus $K_x\ne0$ is equivalent to
$x\in\PGL_2(\Q)\cdot\vartheta$ and to a rational linear dependence among
$1,\vartheta,x,x\vartheta$.  The external prime-valuation vectors of $M$ and $A$ then give a
four-way classification: a $\{2,3\}$-smooth $M$; a $\{2,3\}$-smooth $A$; nonsmooth outputs with a common
external valuation ray; or two independent external valuation directions.

Finally, if $\alpha\in\Apos\setminus\{1\}$, then at least one of
$\alpha^x,\alpha^{x^2},\alpha^{x^3}$ is transcendental, and the logarithm of the first
transcendental term lies outside $\LL$.  The exact failure depth is $1$, $2$, or $3$
according as $\alpha\notin\Gamma$, $\alpha\in\Gamma\setminus K_x$, or
$\alpha\in K_x\setminus\{1\}$.  This yields, for example, the new equivalence that
Alaoglu--Erd\H{o}s holds if and only if every two-base solution makes both
$2^{x^2}$ and $2^{x^3}$ algebraic.  The report closes with a branch-sensitive research
program, emphasizing a two-prime tropical/$p$-adic determinant in the generic branch and
pure-radical or common-core methods in the exceptional branches.
\end{abstract}

\tableofcontents
\clearpage
\pagenumbering{arabic}

\section{Executive answer and trust boundaries}

\subsection{What was and was not proved}

The requested primary goal was an unconditional proof.  After reconstructing the present
formal and paper-level frontier, trying direct logarithmic, determinant, radical, valuation,
and orbit arguments, and then pushing the auxiliary-base and iterated-output structures, I do
not have such a proof.  I found no valid step that turns the exact relation
\[
  \log M\,\log 3=\log A\,\log 2,
  \qquad M=2^x\in\N,\quad A=3^x\in\N,
\]
into a forbidden linear form in logarithms.  This relation is bilinear, and its normalized
$2\times2$ logarithm matrix is precisely the unresolved four-exponentials configuration.
The open status of the general four exponentials conjecture was rechecked against a reviewed
July 2026 problem record and the standard transcendence literature~\cite{TheoremDBFourExp,Waldschmidt2003,Waldschmidt2005}.

That negative verdict should not obscure the amount of genuine progress.  The report contains
three kinds of statements:

\begin{itemize}
\item \statuslean{} statements already checked by the Lean kernel in the repository;
\item \statuspaper{} statements proved in the three inspected LaTeX reports but not necessarily
      formalized;
\item \statusnew{} statements proved in this report and independently checked at the paper level,
      but not yet machine-checked or peer-reviewed.
\end{itemize}

The computational exclusion through $2^{26}$ is marked \statuscomp.  It uses directed MPFR
intervals and a complete finite scan according to the supplied source, but it is not a Lean
proof and I did not rerun all $67{,}108{,}863$ nonzero inputs in this environment.  I did audit
the rounding logic, endpoint logic, and nearest-integer coverage described by the source.

\subsection{The strongest additions in this report}

The following are the main additions relative to the inspected article corpus.

\begin{longtable}{@{}>{\raggedright\arraybackslash}p{0.27\textwidth}>{\raggedright\arraybackslash}p{0.655\textwidth}@{}}
\toprule
Result & Content \\
\midrule
\endhead
Saturated-group synthesis & The current Lean divisible-hull theorem is recast as the
$2,3$ case of a general two-base spectrum
$\GammaSat(a,b)=\{a^rb^s:r,s\in\Q\}$, with a short six-exponential proof. \\
Strong spectrum & Outside $\Gamma$, not only is $\alpha^x$ transcendental, but
$x\log\alpha\notin\LL$ by Baker plus Roy. \\
Valuation lattice & Within $\Gamma$, rational and integral outputs are classified exactly by
the prime-valuation matrix of $(M,A)$; Smith normal form gives the finite rational-output
defect. \\
Symmetric constants & For $\vartheta=\log_2 3$, at least one of
$3^{\vartheta}$ and $2^{1/\vartheta}$ has logarithm outside $\LL$; under a counterexample,
each is algebraic exactly in one explicit smooth-output branch. \\
Second-iterate kernel & The set of $2^r3^s$ whose $x$-th powers return to $\Gamma$ is a
solution-independent vector space of dimension at most one. \\
M\"obius reformulation & The kernel is nonzero exactly when
$x\in\PGL_2(\Q)\cdot\vartheta$, equivalently when $1,\vartheta,x,x\vartheta$ are rationally
dependent. \\
Depth-three tower theorem & Every nontrivial positive algebraic base has a transcendental term
among $\alpha^x,\alpha^{x^2},\alpha^{x^3}$; the first failing depth is classified exactly. \\
New conjecture equivalences & Any algebraic base outside $\Gamma$ is a one-output detector;
for bases inside $\Gamma$, three tower levels are enough.  In particular, base $2$ alone at
levels $x^2$ and $x^3$ detects the conjecture. \\
\bottomrule
\end{longtable}

The exact $2,3$ algebraic-output locus itself is \statuslean{} in the current snapshot;
this report's new content begins with its general saturated-group packaging, the Roy-strengthened
form, and the subsequent valuation-kernel consequences.  Here ``new'' means not isolated in
the inspected article corpus and derived during this analysis.  No claim of priority in the
broader literature is intended.

\subsection{Repository snapshot}

The audit used commit
\[
\sha{4848c6d834190301e014028a9c78a3f0ef49b7b8},
\]
whose root toolchain selects Lean~$4.32.0$.  The article glob contains three source reports:

\begin{enumerate}
\item \texttt{alaoglu-erdos-combined-report/alaoglu-erdos-combined-report.tex};
\item \texttt{alaoglu-erdos-further-report/alaoglu-erdos-further-report.tex};
\item \texttt{alaoglu-erdos-report-3/alaoglu-erdos-report-3.tex}.
\end{enumerate}

They are cumulative rather than independent~\cite{ProveItCombined,ProveItFurther,ProveItReport3}: the combined report consolidates the original
formal program and, at this snapshot, includes the generic positive-real algebraic
determinant, the exact arbitrary positive-algebraic output locus, and second-iterate rank
constraints; the further report develops arbitrary-node and mixed-exponent interpolation
determinants; the third report pushes nonlinear closure, primitive-solution statistics,
finite computation, and mixed-output transcendence.

\section{The problem and its exact transcendence-theoretic shape}

\subsection{Elementary reductions}

Let
\[
 \Sol=\{x\in\R:2^x\in\Z,\ 3^x\in\Z\}.
\]
Because both powers are positive, their integral values are positive.  If $x<0$, then
$0<2^x<1$, so $x\notin\Sol$.  Hence $\Sol\subseteq[0,\infty)$.

\begin{lemma}[Rational solutions are integral]\label{lem:rational-integral}
If $x\in\Q$ and $2^x\in\Z$, then $x\in\Z$.  Consequently every nonintegral element of
$\Sol$ is irrational.
\end{lemma}

\begin{proof}
Write $x=p/q$ in lowest terms with $q>0$.  If $2^{p/q}=M\in\N$, then $2^p=M^q$.
Unique factorization gives $q\mid p$, hence $q=1$.
\end{proof}

If $x$ were algebraic irrational, the Gelfond--Schneider theorem~\cite{Gelfond1960,Lang1966} would make $2^x$
transcendental.  Therefore:

\begin{proposition}[Transcendence of a counterexample]\label{prop:x-trans}
Every nonintegral $x\in\Sol$ is transcendental.
\end{proposition}

This is both elementary from Gelfond--Schneider and formalized in the repository in the
appropriate real-power language.

\subsection{Candidate integers and the slope \texorpdfstring{$\vartheta$}{vartheta}}

Set
\[
  \vartheta=\frac{\log 3}{\log 2}=\log_2 3.
\]
Then $\vartheta$ is irrational by unique factorization and transcendental by
Gelfond--Schneider~\cite{Gelfond1960}, since $2^{\vartheta}=3$.  If $M=2^x$, then
\[
  3^x=M^{\vartheta}.
\]
Thus the original problem is equivalent to classifying the candidate set
\[
 \Cand=\{M\in\N_{>0}:M^{\vartheta}\in\N\}.
\]
The conjecture is
\[
  \Cand=\{2^n:n\in\Nzero\}.
\]
The logarithmic relation for a candidate-output pair $(M,A)$ is
\[
 \frac{\log M}{\log 2}=\frac{\log A}{\log 3}=x.
\]
Equivalently,
\[
 \det\begin{pmatrix}
   \log 2&\log 3\\
   \log M&\log A
 \end{pmatrix}=0.                                      \tag{2.1}\label{eq:four-wall}
\]
Under a nonintegral solution, both rows and both columns are linearly independent over
$\Q$.  All four entries are logarithms of algebraic numbers.  The four exponentials
conjecture predicts that such a singular matrix cannot exist~\cite{Waldschmidt2003,Waldschmidt2005}.  This is the exact global wall.

\subsection{Why integrality still matters}

Although \eqref{eq:four-wall} identifies an open transcendence core, the hypotheses are much
stronger than arbitrary algebraicity: $M$ and $A$ are positive integers, and every
$n+kx$ is again a solution.  This supplies:

\begin{itemize}
\item a discrete additive semigroup of exponents;
\item two synchronized multiplicative semigroups of integer outputs;
\item prime valuations and perfect-power descent;
\item exact denominator control after reducing modulo integer translations;
\item interpolation determinants with integral entries;
\item a dense irrational rotation after taking fractional parts.
\end{itemize}

The repository's main achievement is to exploit nearly all of this structure without
silently assuming the missing four-exponentials step.

\section{Audit of the current ProveIt frontier}

\subsection{Kernel-verified backbone}

The aggregate Lean file imports a broad collection of modules under
\texttt{TwoBaseIntegerExponent}.  The most important groups are as follows.

\begin{longtable}{@{}>{\raggedright\arraybackslash}p{0.25\textwidth}>{\raggedright\arraybackslash}p{0.65\textwidth}@{}}
\toprule
Theme & Representative formal content \\
\midrule
\endhead
Finite exclusion & Exact checks below $256$, $4096$, and related localized intervals; in
particular a nonintegral solution has $x\ge 12$. \\
Rationality and transcendence & Rational solutions are integral; nonintegral solutions and
$\log 3/\log 2$ are transcendental. \\
Localization and radicals & Odd-core reduction, canonical radical form, radical degree,
$3$-denominator normalization, and valuation sharpness. \\
Progressions and density & Sharp progression obstructions, higher differences, zero density,
noninteger density, shrinking-target and denominator-growth estimates. \\
Conditional structure & Least solution, primitive generator, exact counting, quadratic
asymptotics, block Weyl statements, output normalization, common-power rigidity. \\
Auxiliary outputs & Rational and algebraic third bases, generic positive-real algebraic
six-exponentials, exact rational- and algebraic-base output loci, and rank-three
reformulations. \\
Second iterates and external cores & Mixed-output half-plane transcendence, rank-one ratio
loci and an explicit external-prime valuation line, core independence, and canonical-radical
divisible-hull exclusions. \\
\bottomrule
\end{longtable}

Three current formal interfaces are especially important for this report.  First, for every
positive algebraic real $\gamma$ under the two integrality hypotheses,
\[
 \gamma^x\in\Qbar
 \quad\Longleftrightarrow\quad
 x\in\Z\ \text{or}\quad
 \exists n\ge1,\;\gamma^n\in\langle2,3\rangle_{\Z}.
\]
For a nonintegral solution, the right-hand exceptional locus is exactly
$\Gamma=\{2^r3^s:r,s\in\Q\}$.  This is formalized through
\lean{HasPositiveTwoThreeUnitPower} and
\lean{TwoBaseNonintegerSolution.real_rpow_isAlgebraic_iff_hasPositiveTwoThreeUnitPower}.
The natural-base classification
\[
 a^x\in\Qbar
 \quad\Longleftrightarrow\quad
 x\in\Z\ \text{or}\ a=2^u3^v\quad(u,v\in\Nzero)
\]
is a special case.

Second, the new \lean{AlgebraicOutputLocus} module proves that every genuinely mixed base
$(2^i3^j)^{x^2}$ with $i,j>0$ is transcendental under a nonintegral solution, and that two
algebraic members of the ratio family $(2^i/3^j)^{x^2}$ must have collinear exponent vectors.
Third, the corresponding equivalences include the base-$5$ algebraicity formulation, the
base-$6$ squared-exponent formulation, the positive-algebraic rank-three locus, and the
rank-two ratio-locus formulation of Alaoglu--Erd\H{o}s.

\subsection{Exact conditional semigroup structure}

The central conditional theorem can be summarized as follows~\cite{ProveItCombined}.

\begin{theorem}[Conditional rank-two structure; corpus synthesis]\label{thm:conditional-monoid}
Assume the conjecture fails.  Then there is a unique least nonintegral solution $\beta>0$,
and
\[
  \Sol=\{n+k\beta:n,k\in\Nzero\}.
\]
Writing
\[
  M=2^{\beta},\qquad A=3^{\beta},
\]
one has exact output formulae
\[
  2^{n+k\beta}=2^nM^k,
  \qquad
  3^{n+k\beta}=3^nA^k,
\]
and therefore
\[
  \Cand=\{2^nM^k:n,k\in\Nzero\}.
\]
The representation is unique.
\end{theorem}

The existence of a least solution uses discreteness forced by the candidate integers;
closure under subtraction of ordered solutions then supplies Euclidean descent.  Uniqueness
of the $(n,k)$ representation uses the transcendence of $\beta$.  This is considerably more
rigid than a generic hypothetical set of exceptional exponents.

The least pair $(M,A)$ is affinely primitive: one cannot simultaneously remove a common
integer translation from the exponent or a common perfect-power factor from both outputs.
The current arithmetic-rigidity module formalizes a quantitative version: if
\[
  2^x=2^r a^k,\qquad 3^x=3^r b^k,
\]
then $(x-r)/k$ is another solution; the verified zero-free interval yields
\[
  r+8k\le x
\]
for a nonintegral solution.  Later paper-level finite work improves the raw exclusion but the
formal affine theorem deliberately depends only on the kernel-verified bound available in
that module.

\subsection{Counting, spacing, and interpolation determinants}

Conditional on failure, exact rank-two structure gives logarithmic local counts and quadratic
cumulative counts.  The further report also proves unconditional candidate repulsion from
arbitrary-node generalized Vandermonde determinants~\cite{ProveItFurther}.  A typical statement is that if
\[
  N\le m_0<m_1<\cdots<m_r\le N+H,
  \qquad m_j\in\Cand,
\]
then
\[
  \prod_{0\le i<j\le r}(m_j-m_i)
  \ge
  \frac{r!}{\left|\vartheta(\vartheta-1)\cdots(\vartheta-r+1)\right|}
  m_0^{\,r-\vartheta}.
\]
Using mixed exponents $a+b\vartheta$ strengthens the local spacing scale; for sufficiently
large $r$ the report obtains a bound of the schematic form
\[
  H\ge \frac1{24}N^{1-12/\sqrt r}.
\]
The resulting unconditional dyadic count is $O((\log X)^2)$, while a counterexample's exact
monoid predicts only $O(\log X)$ points in $[X,2X]$.  This single missing logarithm is a useful
diagnosis: the determinant method is close enough to see the correct geometry, but not close
enough to contradict it.

\subsection{Nonlinear closure and the third report}

The third report~\cite{ProveItReport3} exploits the transcendence of $\beta$ inside the exact monoid.  Among its
cleanest paper-level statements are:
\[
  (n+k\beta)(m+\ell\beta)\in\Sol
  \quad\Longleftrightarrow\quad
  k=0\ \text{or}\ \ell=0,
\]
so two nonintegral solutions never multiply to another solution.  More generally, polynomial
and rational-function intersections with $\Sol$ collapse to affine functions in $\beta$.
It also derives primitive-solution asymptotics, exact divisor-lattice formulae, root-closure
criteria, optimized high-difference exclusions, and mixed-exponent transcendence such as the
transcendence of $6^{xy}$ for nonintegral solutions $x,y$.

Some of these statements have since acquired neighboring Lean infrastructure, but the whole
third-report package is not yet represented by one formal module.  This report therefore
keeps its status labels theorem-by-theorem rather than treating the PDF as a machine-checked
object.

\section{Attempts at a complete proof and the precise points of failure}

This section records the principal proof attempts.  Several produce useful lemmas later in
the report; none produces the missing contradiction.  It is important to separate a route
that is merely unfinished from one that has silently restated the open conjecture.

\subsection{Direct linear forms in logarithms}

The relation
\[
  \log M\,\log3-\log A\,\log2=0
\]
is not a linear form in logarithms.  Baker's theorem controls expressions
\[
  c_1\log\alpha_1+\cdots+c_n\log\alpha_n
\]
with algebraic coefficients $c_i$; here the coefficients $\log 2$ and $\log 3$ are themselves
transcendental logarithms.  Dividing by $\log2\log3$ merely returns
\[
  \frac{\log M}{\log2}=\frac{\log A}{\log3}.
\]
Any argument saying that the equality of these logarithmic ratios forces multiplicative
dependence of $(2,M)$ or $(3,A)$ is exactly a four-exponentials assertion and is not presently
available.

The integrality of $M,A$ does permit prime expansions
\[
 \log M=\sum_p \vpp pM\log p,
 \qquad
 \log A=\sum_p \vpp pA\log p,
\]
but substitution yields a quadratic relation among logarithms:
\[
 \sum_p \vpp pM\log p\log3
 -\sum_p \vpp pA\log p\log2=0.
\]
Known homogeneous linear-independence results do not control such a relation.  The support
information will become useful only after an additional algebraic output is manufactured.

\subsection{Trying to manufacture a third algebraic output}

The six exponentials theorem~\cite{Ramachandra1968,Lang1966,Waldschmidt2000} immediately finishes the problem if one can find a positive
algebraic base $\alpha$, multiplicatively independent of $2$ and $3$, such that $\alpha^x$ is
algebraic.  Indeed, using rows $\log2,\log3,\log\alpha$ and columns $1,x$ would give six
algebraic exponentials, contradicting six exponentials.

The current Lean module \texttt{AlgebraicThirdBase.lean} realizes this principle internally
for natural bases by an interpolation determinant and upgrades rationality to integrality.
Its consequence is not a proof of Alaoglu--Erd\H{o}s; it says that a counterexample forces
\emph{every} independent natural auxiliary output to be transcendental.

Natural candidates for a third base include $M+A$, $M-A$ when positive, prime divisors of
$M$ or $A$, and primitive odd cores.  They are algebraic bases, but no identity makes their
$x$-th powers algebraic.  The exact spectrum theorem in the next section shows that this is
not merely a failure of imagination: under a counterexample, no algebraic base outside the
saturated group generated by $2$ and $3$ can have algebraic output.  Thus a successful
``third-base'' proof must derive algebraicity from an additional arithmetic operation, not
just choose a clever base.

\subsection{Radical normalization and local ramification}

A candidate $M$ can be written in a canonical radical form.  Removing the largest power of
$2$, extracting the maximal perfect power from the odd core, and normalizing powers of $3$
reduces the conjecture to a statement of the form
\[
 u^{\vartheta}\notin\Z[1/3]
 \qquad\text{for every odd power-primitive }u>1
\]
subject to explicit side conditions.  The repository proves exact radical-degree and
valuation statements around this reduction.

A tempting route is to place a hypothetical value
\[
  u^{\vartheta}=\frac{a}{3^t}
\]
in a pure extension and compare ramification at $3$ and at primes dividing $u$.  This does
not work directly because $\vartheta$ is transcendental: the notation is an equality of
positive real numbers, not an algebraic radical extension generated by a rational exponent.
One can obtain genuine radicals after a smooth-output or common-perfect-power reduction, but
then the remaining relation is again a product of logarithms.  Local valuations rigorously
control every \emph{algebraic} power that appears; they do not assign a $p$-adic meaning to the
real exponent $\vartheta$ itself.

A productive reformulation is therefore branch-sensitive.  In the smooth-output branches
of \cref{sec:symmetric}, one of the constants $3^{\vartheta}$ or $2^{1/\vartheta}$ becomes a
genuine algebraic radical, and local Kummer theory can legitimately be brought to bear.  In
the generic branch, two independent external valuation directions are available instead.

\subsection{Integer-gap estimates from rational approximation}

Let $p/q$ approximate $x$.  Then
\[
 M^q-2^p\in\Z,
 \qquad
 A^q-3^p\in\Z.
\]
If neither difference vanishes, each has absolute value at least $1$.  Taylor expansion gives
bounds of the rough form
\[
  |qx-p|\gg 2^{-p},
  \qquad
  |qx-p|\gg 3^{-p}.
\]
These are valid denominator-growth estimates but exponentially weaker than the polynomial
lower bounds that would conflict with continued fractions.  The denominator of the rational
approximation is $q$, whereas the integer scale is exponential in $q$.

One can cancel several Taylor terms with Hermite--Pad\'e or interpolation determinants.  The
further report does exactly this in a coordinate-free way and obtains strong local repulsion,
but the exact counterexample monoid remains sparse enough to fit within the bounds.  A cluster
built from a convergent $p/q$ to the least solution $\beta$ has baseline size exponential in
$q$; its small relative width does not become a sufficiently small \emph{absolute} integer
quantity.

\subsection{Archimedean interpolation determinants}

For candidate pairs $(m_i,n_i)$ with $n_i=m_i^{\vartheta}\in\N$, matrices with entries
\[
  m_i^{a_j}n_i^{b_j}=m_i^{a_j+b_j\vartheta}
\]
have integer determinants.  If the nodes $m_i$ are close, confluent Vandermonde estimates make
the determinant small in the ordinary absolute value.  Nonvanishing then gives a lower bound
of $1$, producing spacing inequalities.

There are three recurring limitations.

\begin{enumerate}
\item The number of useful mixed exponents grows quadratically with the exponent box, while
      the analytic losses also grow with determinant size.
\item A generic dyadic block estimate sees all possible candidate configurations, not the
      exact triangular $(n,k)$ lattice forced by a counterexample.
\item The archimedean norm alone discards prime-valuation information in the entries.
\end{enumerate}

The first two explain the one-logarithm deficit between the unconditional $O((\log X)^2)$
dyadic bound and the conditional $O(\log X)$ count.  The third motivates the proposed
adelic/tropical determinant in \cref{sec:future}.

\subsection{\texorpdfstring{$p$-adic}{p-adic} exponentiation and why it cannot simply replace the real exponent}

For a $p$-adic unit close to $1$, the function $z\mapsto u^z$ is $p$-adically analytic on a
suitable domain.  However, the real number $x$ has not been shown to define a compatible
$p$-adic exponent, and an integer such as $M=2^x$ need not determine such a compatibility.
Writing the real relation inside a $p$-adic logarithm without first producing an algebraic
identity is invalid.

What \emph{is} legitimate is to apply $p$-adic valuations to the integer interpolation
matrices, to algebraic radicals already certified by rational exponents, or to multiplicative
relations such as $M^rA^s=2^u3^v$.  The second-iterate kernel introduced later cleanly detects
when the latter relations occur.

\subsection{Additive and multiplicative closure}

The solution set is additively closed, while the third report proves that, under failure, the
product of two nonintegral solutions is never a solution.  One might hope to derive an
independent closure under squaring or multiplication from the original power identities.
No such identity is available:
\[
  2^{x^2}=(2^x)^x=M^x
\]
raises an integer base to the same unknown exponent and is generally transcendental.  The
repository's theorem that $6^{x^2}$ is transcendental shows that nonlinear closure fails in a
particularly robust way.  The tower theorem in \cref{sec:tower} explains the exact depth at
which every algebraic base must fail.

\subsection{Bottom line of the proof attempts}

A complete proof would follow from any one of the following genuinely new bridges:

\begin{itemize}
\item algebraicity of $\alpha^x$ for one algebraic $\alpha\notin\Gamma$;
\item a sub-$O(\log X)$ unconditional dyadic count for candidates;
\item a local/$p$-adic obstruction to every one of the four kernel branches in
      \cref{sec:kernel-types};
\item a proof that one of $2^{x^2}$ and $2^{x^3}$ is algebraic under the original hypotheses
      in a way that contradicts the depth-three theorem;
\item the relevant special case of four exponentials, strong four exponentials, or a
      sufficiently strong algebraic-independence conjecture for logarithms.
\end{itemize}

The rest of this report turns the failed third-base and iterated-output attempts into exact
classification theorems.

\section{Exact algebraic auxiliary-base spectrum}\label{sec:spectrum}

\subsection{The saturated algebraic group}

The multiplicative group $\Apos$ becomes a vector space over $\Q$ when rational scalar
multiplication is interpreted by the unique positive real root:
\[
 q\cdot\alpha=\alpha^q,
 \qquad q\in\Q,
 \quad \alpha\in\Apos.
\]
For multiplicatively independent $a,b\in\Apos$, define their saturated group
\[
  \GammaSat(a,b)=\{a^r b^s:r,s\in\Q\}.
\]
The map $\Q^2\to\GammaSat(a,b)$ is an isomorphism of $\Q$-vector spaces.

The current Lean snapshot already contains the essential determinant in a fully symmetric
form: three positive algebraic bases with injective nonnegative monomials cannot all have
algebraic $x$-th powers for an irrational $x$.  It also packages the specialization to
$2,3$ as an exact divisible-hull classification.  The next theorem is the concise
coordinate-free synthesis for an arbitrary independent pair $a,b$; its proof is a direct
application of the same classical six-exponentials input.

\begin{theorem}[Exact algebraic-base spectrum]\label{thm:general-spectrum}
Let $a,b\in\Apos$ be multiplicatively independent, let $x\in\R\setminus\Q$, and assume
$a^x,b^x\in\Qbar$.  Then, for every $\alpha\in\Apos$,
\[
  \alpha^x\in\Qbar
  \quad\Longleftrightarrow\quad
  \alpha\in\GammaSat(a,b).
\]
\statusnew
\end{theorem}

\begin{proof}
If $\alpha=a^rb^s$ with $r,s\in\Q$, positivity gives
\[
 \alpha^x=(a^x)^r(b^x)^s,
\]
which is algebraic.

Conversely, suppose $\alpha\notin\GammaSat(a,b)$.  Then
\[
 \log a,\quad \log b,\quad \log\alpha
\]
are linearly independent over $\Q$.  Indeed, a rational relation can be cleared to an integer
multiplicative relation.  If the coefficient of $\log\alpha$ is zero, it contradicts the
multiplicative independence of $a,b$; otherwise it places $\alpha$ in the saturated group.
The pair $1,x$ is linearly independent over $\Q$.  Apply the six exponentials theorem to the
three rows $\log a,\log b,\log\alpha$ and the two columns $1,x$.  Five of the six exponential
values
\[
  a,\ b,\ \alpha,\ a^x,\ b^x
\]
are algebraic.  Therefore $\alpha^x$ is transcendental.
\end{proof}

This theorem is conceptually stronger than a list of auxiliary bases: it says that two
algebraic outputs have already exhausted the maximum algebraic multiplicative rank allowed by
six exponentials.

\subsection{Specialization to Alaoglu--Erd\H{o}s}

For the original problem write
\[
  \Gamma=\GammaSat(2,3)=\{2^r3^s:r,s\in\Q\}.
\]

\begin{corollary}[Complete algebraic spectrum of a counterexample]\label{cor:AE-spectrum}
Let $x\in\Sol\setminus\Z$.  For every $\alpha\in\Apos$,
\[
  \alpha^x\in\Qbar
  \quad\Longleftrightarrow\quad
  \alpha\in\Gamma.
\]
In particular, the spectrum is independent of the hypothetical nonintegral solution $x$.
\statuslean
\end{corollary}

\begin{proof}
By \cref{lem:rational-integral}, $x$ is irrational.  Applying
\cref{thm:general-spectrum} with $a=2$, $b=3$ proves the statement directly.
Equivalently, the current Lean predicate
\lean{HasPositiveTwoThreeUnitPower} says that some positive natural power of $\alpha$ is a
positive rational $2,3$-unit.  For positive algebraic $\alpha$, this is equivalent to
$\alpha\in\Gamma$: one implication raises $2^r3^s$ to a common denominator, and the other
uses the unique positive real root.  The theorem
\lean{TwoBaseNonintegerSolution.real_rpow_isAlgebraic_iff_hasPositiveTwoThreeUnitPower}
therefore gives exactly the displayed equivalence.
\end{proof}

The group $\Gamma$ is countable, divisible, of rational rank two, and dense in
$\R_{>0}$; density already follows from the subgroup $\{2^r:r\in\Q\}$.  Thus a counterexample
would have a dense set of algebraic bases with algebraic output, but every algebraic base
outside that exact dense subgroup would have transcendental output.

\begin{corollary}[Rank bound for any family of algebraic outputs]
Let $x\in\Sol\setminus\Z$, and let $\alpha_1,\ldots,\alpha_t\in\Apos$ satisfy
$\alpha_i^x\in\Qbar$.  Then the multiplicative rank of the $\alpha_i$ is at most two, and
every $\alpha_i$ lies in $\Gamma$.
\end{corollary}

\subsection{Every independent algebraic base is a detector}

\begin{theorem}[Algebraic detector equivalence]\label{thm:detector}
Fix $\alpha\in\Apos\setminus\Gamma$.  Then the Alaoglu--Erd\H{o}s conjecture is equivalent to
\[
  \forall x\in\R,\quad
  2^x,3^x\in\Z
  \Longrightarrow
  \alpha^x\in\Qbar.
\]
\statusnew
\end{theorem}

\begin{proof}
If the conjecture holds, every solution $x$ is a nonnegative integer, so $\alpha^x$ is
algebraic.  Conversely, a nonintegral solution would make $\alpha^x$ transcendental by
\cref{cor:AE-spectrum}.
\end{proof}

This is a paper-level universalization of the kernel-verified pointwise locus, rather than a
new transcendence estimate.  The repository's base-$5$ equivalence is the natural-integer
instance $\alpha=5$.  The theorem
shows that no arithmetic feature specific to $5$ is needed.  Valid detectors include, for
example,
\[
  5,\quad \sqrt5,\quad 1+\sqrt2,\quad \frac{7+3\sqrt5}{2},
\]
provided the selected number is not accidentally in $\Gamma$.  For an algebraic integer
having a prime divisor outside $2,3$, exclusion from $\Gamma$ is immediate; for a general
algebraic number it may be checked by multiplicative independence.

\begin{corollary}[Zero-one law across nonintegral solutions]
Assume the conjecture fails, and let $\alpha\in\Apos$.  Exactly one of the following holds:
\begin{enumerate}
\item $\alpha\in\Gamma$, and $\alpha^x$ is algebraic for every nonintegral solution $x$;
\item $\alpha\notin\Gamma$, and $\alpha^x$ is transcendental for every nonintegral solution
      $x$.
\end{enumerate}
\end{corollary}

\subsection{Strong spectrum via Baker and Roy}

Let $\LL$ denote the $\Qbar$-vector space generated by $1$ and all logarithms of nonzero
algebraic numbers, allowing arbitrary complex branches.  For positive real algebraic numbers
we use the real logarithm, which is one of those branches.

Roy's strong six exponentials theorem~\cite{Roy1992,Waldschmidt2005} says that if three complex numbers are linearly
independent over $\Qbar$ and two more are linearly independent over $\Qbar$, then at least one
of the six pairwise products lies outside $\LL$.

\begin{theorem}[Strong algebraic-base spectrum]\label{thm:strong-spectrum}
Under the hypotheses of \cref{thm:general-spectrum}, if
$\alpha\notin\GammaSat(a,b)$, then
\[
  x\log\alpha\notin\LL.
\]
In particular, $\alpha^x$ is transcendental.  \statusnew
\end{theorem}

\begin{proof}
As in \cref{thm:general-spectrum}, the three logarithms
$\log a,\log b,\log\alpha$ are $\Q$-linearly independent.  Baker's homogeneous theorem~\cite{Baker1975}
upgrades them to linear independence over $\Qbar$.

The number $x$ is transcendental.  Indeed, if it were algebraic, its irrationality together
with $a\ne0,1$ would make $a^x$ transcendental by Gelfond--Schneider.  Hence $1,x$ are linearly
independent over $\Qbar$.

Apply Roy's theorem to rows $\log a,\log b,\log\alpha$ and columns $1,x$.  Five products lie in
$\LL$:
\[
 \log a,\ \log b,\ \log\alpha,
 \quad x\log a=\log(a^x),
 \quad x\log b=\log(b^x).
\]
Therefore the sixth product $x\log\alpha$ lies outside $\LL$.
\end{proof}

\begin{definition}
For brevity, this report calls a positive exponential value $e^z$ \emph{log-strongly
transcendental} when $z\notin\LL$.  This implies ordinary transcendence of $e^z$, because the
logarithm of a positive algebraic number belongs to $\LL$.
\end{definition}

\begin{corollary}[Strong detector]
For every fixed $\alpha\in\Apos\setminus\Gamma$, Alaoglu--Erd\H{o}s is equivalent to the
assertion that all two-base solutions satisfy
\[
  x\log\alpha\in\LL.
\]
\end{corollary}

\subsection{Five-exponential rays forced by a counterexample}

The classical five exponentials theorem~\cite{Waldschmidt2005} gives another useful description of the hypothetical
world.

\begin{proposition}[Two forbidden exponential rays]\label{prop:five-rays}
Let $x\in\Sol\setminus\Z$, let $\vartheta=\log_2 3$, and let
$0\ne\gamma\in\Qbar$.  Then both
\[
  e^{\gamma x}
  \qquad\text{and}\qquad
  e^{\gamma\vartheta}
\]
are transcendental.  \statusclassical
\end{proposition}

\begin{proof}
For the first value, apply five exponentials to the independent pairs $(1,x)$ and
$(\log2,\log3)$.  The four rectangular exponentials are $2,3,M,A$, all algebraic, so the fifth
$e^{\gamma x}$ is transcendental.

For the second value, use the independent pairs $(\log2,\log3)$ and $(1,x)$.  The same four
rectangular exponentials occur, and the fifth is
$e^{\gamma\log3/\log2}=e^{\gamma\vartheta}$.
\end{proof}

This does not prove the conjecture: $\gamma$ must be algebraic, whereas the useful
coefficients arising from new integer bases are themselves logarithms.  It does show that a
counterexample forces two entire algebraic rays of exponential values to be transcendental.

\section{Rational and integral outputs inside the algebraic spectrum}\label{sec:valuation-spectrum}

The current Lean snapshot settles algebraicity throughout $\Gamma$ (and excludes it
outside $\Gamma$).  A finer arithmetic question remains: if
$\alpha=2^r3^s\in\Gamma$, when is $\alpha^x$ rational or an integer?  For a fixed solution
the answer is completely encoded by prime valuations, and the resulting Smith lattice is a
strict refinement of the divisible-hull theorem.

\subsection{The valuation matrix}

Let $x\in\Sol\setminus\Z$ and write
\[
  M=2^x,\qquad A=3^x.
\]
Let $P$ be the finite union of the prime supports of $M$ and $A$, and define column vectors
\[
  \mathbf m=(\vpp pM)_{p\in P},
  \qquad
  \mathbf a=(\vpp pA)_{p\in P}.
\]
Let
\[
  V_x=\begin{pmatrix}\mathbf m&\mathbf a\end{pmatrix}
  \in M_{P\times2}(\Z).
\]
For $q=(r,s)^T\in\Q^2$ put
\[
  \alpha(q)=2^r3^s.
\]
Then the unique positive rational-power interpretation gives
\[
  \alpha(q)^x=M^rA^s
  =\prod_{p\in P}p^{r\vpp pM+s\vpp pA}.               \tag{6.1}\label{eq:valuation-output}
\]

\begin{lemma}[Rank two]\label{lem:V-rank-two}
The two columns of $V_x$ are linearly independent over $\Q$.
\end{lemma}

\begin{proof}
If $r\mathbf m+s\mathbf a=0$, then $M^rA^s=1$.  Taking real logarithms gives
\[
  x(r\log2+s\log3)=0.
\]
Since $x>0$ and $\log3/\log2$ is irrational, $r=s=0$.
\end{proof}

\begin{theorem}[Exact rational and integral auxiliary-output spectrum]
\label{thm:rational-integral-spectrum}
Let $x,M,A,V_x$ be as above and let $q=(r,s)^T\in\Q^2$.  Then
\begin{align*}
 \alpha(q)^x\in\Q_{>0}
 &\quad\Longleftrightarrow\quad V_xq\in\Z^P,\\
 \alpha(q)^x\in\N
 &\quad\Longleftrightarrow\quad V_xq\in\Nzero^P.
\end{align*}
Together with \cref{cor:AE-spectrum}, this classifies algebraic, rational, and integral
outputs for every positive algebraic auxiliary base.  \statusnew
\end{theorem}

\begin{proof}
Equation \eqref{eq:valuation-output} gives a finite product of primes with rational exponents.
Choose a common denominator $d$.  If this product is rational, raising to the $d$-th power and
comparing ordinary rational prime valuations shows that every exponent in
\eqref{eq:valuation-output} is an integer.  Conversely integer exponents give a rational
number.  Such a rational number is a positive integer exactly when all its prime valuations
are nonnegative.
\end{proof}

Thus three nested spectra occur:
\[
 \begin{array}{rcl}
 \mathcal B_{\rm alg}(x)
  &=&\{\alpha(q):q\in\Q^2\}=\Gamma,\\[2mm]
 \mathcal B_{\rm rat}(x)
  &=&\{\alpha(q):q\in\Lambda_x\},\\[1mm]
 \mathcal B_{\rm int}(x)
  &=&\{\alpha(q):q\in\Lambda_x^+\},
 \end{array}
\]
where
\[
 \Lambda_x=\{q\in\Q^2:V_xq\in\Z^P\},
 \qquad
 \Lambda_x^+=\{q\in\Lambda_x:V_xq\ge0\}.
\]

\subsection{Smith normal form and the finite rational-output defect}

Because $V_x$ has rank two, $\Lambda_x$ is a rank-two lattice in $\Q^2$ containing $\Z^2$
with finite quotient.  Let
\[
 \delta_1(V_x)=\gcd\{\text{all entries of }V_x\},
\]
and let
\[
 \delta_2(V_x)=\gcd\left\{
  \left|\det\begin{pmatrix}
   \vpp pM&\vpp pA\\
   \vpp qM&\vpp qA
  \end{pmatrix}\right|:p,q\in P
 \right\}.
\]
The second gcd is positive by \cref{lem:V-rank-two}.

\begin{theorem}[Smith classification]\label{thm:smith-lattice}
Let $d_1,d_2$ be the Smith invariant factors of $V_x$, so
$d_1=\delta_1(V_x)$, $d_1\mid d_2$, and $d_1d_2=\delta_2(V_x)$.  After a unimodular change of
coordinates in $\Q^2$,
\[
  \Lambda_x=\frac1{d_1}\Z\oplus\frac1{d_2}\Z.
\]
Consequently
\[
  \Lambda_x/\Z^2\cong\Z/d_1\Z\oplus\Z/d_2\Z,
  \qquad
  |\Lambda_x/\Z^2|=\delta_2(V_x).
\]
\statusnew
\end{theorem}

\begin{proof}
Choose unimodular matrices $U,W$ with
\[
 UV_xW=\begin{pmatrix}d_1&0\\0&d_2\\0&0\\\vdots&\vdots\end{pmatrix}.
\]
Since $U$ preserves the integer lattice in the target and $W$ preserves $\Z^2$ in the
domain, the condition $V_xq\in\Z^P$ becomes
$d_1q'_1,d_2q'_2\in\Z$ in the transformed coordinates.  The quotient and determinant-divisor
formulae are the standard Smith-normal-form identities.
\end{proof}

The first invariant has a direct output interpretation: $d_1$ is the largest integer $d$
for which both $M$ and $A$ are perfect $d$-th powers.  For the least conditional solution
$\beta$, affine primitivity gives $d_1=1$.  In that case
\[
  \Lambda_\beta/\Z^2\cong\Z/\delta_2(V_\beta)\Z.
\]
Thus one integer, the gcd of all $2\times2$ valuation minors, measures the entire finite defect
between rational $\beta$-outputs at algebraic $2,3$-radical bases and the obvious rational
$2,3$-unit bases.

The integral cone $\Lambda_x^+$ is the intersection of this lattice with the rational
polyhedral cone $V_xq\ge0$.  By Gordan's lemma it is a finitely generated normal affine
semigroup.  This is a useful exact replacement for trying to enumerate algebraic auxiliary
bases one at a time.

\subsection{Two translated solutions remove the fractional defect}

The algebraic spectrum does not depend on $x$, but the rational-output lattice $\Lambda_x$
can.  Translation by one gives a sharp intersection identity.

\begin{proposition}[Translation intersection]\label{prop:translation-lattice}
Let $x\in\Sol\setminus\Z$.  Then
\[
  \Lambda_x\cap\Lambda_{x+1}=\Z^2.
\]
Equivalently, for $\alpha\in\Apos$, if both $\alpha^x$ and $\alpha^{x+1}$ are rational, then
\[
  \alpha=2^u3^v\qquad(u,v\in\Z).
\]
\statusnew
\end{proposition}

\begin{proof}
By the algebraic spectrum, write $\alpha=2^r3^s$.  The quotient
\[
  \alpha=\frac{\alpha^{x+1}}{\alpha^x}
\]
is rational.  A rational number of the form $2^r3^s$ with $r,s\in\Q$ must have
$r,s\in\Z$, by ordinary prime valuations.  The converse is immediate.
\end{proof}

\begin{corollary}[Universal rational and integral spectra under failure]
Assume the conjecture fails and let $\alpha\in\Apos$.
\begin{enumerate}
\item $\alpha^y$ is rational for every nonintegral solution $y$ if and only if
      $\alpha=2^u3^v$ with $u,v\in\Z$.
\item $\alpha^y$ is an integer for every solution $y$ if and only if
      $\alpha=2^u3^v$ with $u,v\in\Nzero$.
\end{enumerate}
\end{corollary}

\begin{proof}
For the first part, use any nonintegral $x$ and $x+1$ in
\cref{prop:translation-lattice}.  Conversely rational $2,3$-units have rational outputs
$M^uA^v$.

For the second part, the forward implication first gives
$\alpha=a/b\in\Q_{>0}$ in lowest terms.  If $z=\alpha^x\in\N$ and
$z(a/b)^n=\alpha^{x+n}$ is an integer for every $n\ge0$, then $b^n$ divides the fixed integer
$z$ for every $n$, so $b=1$.  An integer in $\Gamma$ is exactly $2^u3^v$ with
$u,v\in\Nzero$.  The converse is immediate.
\end{proof}

\section{The symmetric tower constants}\label{sec:symmetric}

The constants
\[
  \vartheta=\frac{\log3}{\log2},
  \qquad
  \eta=\frac1\vartheta=\frac{\log2}{\log3}
\]
generate two symmetric ``next outputs''
\[
  E_3=3^{\vartheta},
  \qquad
  E_2=2^{\eta}.
\]
Neither individual constant is known to be transcendental.  Six exponentials does, however,
control the pair, and a hypothetical counterexample makes their arithmetic significance
exact.

\subsection{An unconditional strong dichotomy}

\begin{theorem}[At least one symmetric tower constant is strongly transcendental]
\label{thm:symmetric-unconditional}
At least one of
\[
  \vartheta\log3=\log E_3,
  \qquad
  \eta\log2=\log E_2
\]
lies outside $\LL$.  In particular, at least one of $E_3,E_2$ is transcendental.
\statusnew
\end{theorem}

\begin{proof}
The three numbers $1,\vartheta,\eta$ are linearly independent over $\Qbar$.  A relation
$a+b\vartheta+c/\vartheta=0$ would give
$b\vartheta^2+a\vartheta+c=0$ with algebraic coefficients, impossible because $\vartheta$ is
transcendental.  The logarithms $\log2,\log3$ are linearly independent over $\Qbar$ by
Baker's theorem.

Apply Roy's strong six exponentials theorem to those three rows and two columns.  Four of the
six products are already logarithms of algebraic numbers:
\[
 \log2,\quad \log3,\quad
 \vartheta\log2=\log3,\quad
 \eta\log3=\log2.
\]
Therefore one of the two remaining products $\vartheta\log3$ and $\eta\log2$ lies outside
$\LL$.
\end{proof}

The ordinary six exponentials theorem gives the weaker transcendence dichotomy directly.  The
strong form is useful because it persists in the branch classification below.

\subsection{Exact conditional classification of \texorpdfstring{$E_3$}{E3}}

Let $x\in\Sol\setminus\Z$, $M=2^x$, and $A=3^x$.

\begin{theorem}[The $E_3$ equivalences]\label{thm:E3}
The following are equivalent:
\begin{enumerate}
\item $E_3=3^{\vartheta}$ is algebraic;
\item $\vartheta\log3\in\LL$;
\item $M$ is $\{2,3\}$-smooth;
\item there exist $a,b\in\Nzero$ with $b>0$ such that
      \[
        M=2^a3^b,
        \qquad x=a+b\vartheta;
      \]
\item $x\in\Q+\Q\vartheta$.
\end{enumerate}
If these conditions fail, then $\vartheta\log3\notin\LL$ and $E_3$ is log-strongly
transcendental.  If they hold, then
\[
  E_3^b=\frac{A}{3^a}\in\Q_{>0},
\]
so $E_3$ is an explicit algebraic radical.  \statusnew
\end{theorem}

\begin{proof}
The implications (4)$\Rightarrow$(3)$\Rightarrow$(5) are immediate, since taking logarithms
of $M=2^a3^b$ gives $x=a+b\vartheta$.  Conversely, if (5) holds, write
$x=u+v\vartheta$ with $u,v\in\Q$ and clear a common positive denominator $d$:
\[
  M^d=2^p3^q
  \qquad(p,q\in\Z).
\]
Because the left side is a positive integer, prime valuations force $p,q\ge0$ and
$d\mid p,q$.  Thus $u=a$ and $v=b$ for $a,b\in\Nzero$; nonintegrality of $x$ gives
$b>0$.  Hence (5)$\Rightarrow$(4).

Suppose (1) holds.  Apply six exponentials to rows $1,x,\vartheta$ and columns
$\log2,\log3$.  The six exponentials are
\[
  2,\ 3,\ M,\ A,\ 3,\ E_3,
\]
all algebraic.  Hence the three rows are rationally dependent, so
$x=u+v\vartheta$ for $u,v\in\Q$.  The implication (5)$\Rightarrow$(4) just proved now gives
$M=2^a3^b$ with $a,b\in\Nzero$ and $b>0$.  This proves (1)$\Rightarrow$(4), while (4) gives
$A=3^aE_3^b$ and hence (1).

It remains to prove the strong statement.  If $M$ is not smooth, then
$\log2,\log M,\log3$ are linearly independent over $\Q$: a rational relation with a nonzero
$\log M$ coefficient would put $M$ in $\Gamma$, and an integer in $\Gamma$ is smooth.  Baker
upgrades this to $\Qbar$-linear independence.  Therefore $1,x,\vartheta$ are
$\Qbar$-linearly independent after multiplication by $\log2$.  Roy's theorem applies to
rows $1,x,\vartheta$ and columns $\log2,\log3$.  The other five products lie in $\LL$, forcing
$\vartheta\log3\notin\LL$.  Thus failure of (3) implies failure of (2), so (2) implies
(3); conversely, (1) plainly implies (2).
\end{proof}

\subsection{The symmetric \texorpdfstring{$E_2$}{E2} classification}

\begin{theorem}[The $E_2$ equivalences]\label{thm:E2}
The following are equivalent:
\begin{enumerate}
\item $E_2=2^{\eta}$ is algebraic;
\item $\eta\log2\in\LL$;
\item $A$ is $\{2,3\}$-smooth;
\item there exist $c,d\in\Nzero$ with $c>0$ such that
      \[
        A=2^c3^d,
        \qquad x=d+c\eta;
      \]
\item $x\in\Q+\Q\eta$.
\end{enumerate}
If these conditions fail, $\eta\log2\notin\LL$.  If they hold, then
\[
  E_2^c=\frac{M}{2^d}\in\Q_{>0}.
\]
\statusnew
\end{theorem}

\begin{proof}
Interchange $2$ and $3$, and interchange $\vartheta$ with $\eta$, in the proof of
\cref{thm:E3}.  For the strong step, use the $\Q$-independence of
$\log3,\log A,\log2$ when $A$ is not smooth.
\end{proof}

\begin{proposition}[The smooth cases are mutually exclusive]\label{prop:not-both-smooth}
For a nonintegral solution, $M$ and $A$ cannot both be $\{2,3\}$-smooth.  Consequently
$E_3$ and $E_2$ cannot both be algebraic.
\end{proposition}

\begin{proof}
Write
\[
  M=2^a3^b,
  \qquad
  A=2^c3^d.
\]
Then
\[
 x=a+b\vartheta=d+c\eta.
\]
Multiplying the second equality by $\vartheta$ and eliminating $x$ gives
\[
 b\vartheta^2+(a-d)\vartheta-c=0.
\]
The transcendence of $\vartheta$ forces $b=c=0$ and $a=d$, hence $x=a\in\Z$, a
contradiction.
\end{proof}

The current Lean development proves a more refined primitive-core exclusion by deriving a
quadratic polynomial for $\vartheta$.  The argument above is its smooth-output shadow and
will fit naturally into the kernel classification.

\section{The solution-independent second-iterate kernel}\label{sec:kernel}

The exact spectrum controls one application of the exponent $x$.  The next question is which
bases in $\Gamma$ return to $\Gamma$ after that application.  This turns mixed second
iterates into finite-dimensional linear algebra on prime valuations.

The current module \lean{AlgebraicOutputLocus} already proves two substantial slices of this
picture: every genuinely mixed $(2^i3^j)^{x^2}$ with $i,j>0$ is transcendental, and the
algebraic rational-ratio outputs lie on one explicit external $p$-adic line (hence have
lattice rank at most one).  The kernel $K$ below upgrades those nonnegative-lattice statements
to an exact $\Q$-linear invariant, valid simultaneously for every nonintegral solution and
every rational direction.

\subsection{Definition and valuation description}

For $x\in\Sol\setminus\Z$ define
\[
 K_x=\left\{(r,s)\in\Q^2:M^rA^s\in\Gamma\right\}.
                                                        \tag{8.1}\label{eq:K-def}
\]
Equivalently, $q=(r,s)$ lies in $K_x$ exactly when
\[
  \alpha(q)^x\in\Gamma,
  \qquad \alpha(q)=2^r3^s.
\]
Because $q\mapsto\alpha(q)$ is injective, we will also identify $K_x$ with the corresponding
rank-at-most-one subgroup $\{\alpha(q):q\in K_x\}\subseteq\Gamma$ when discussing
power orbits.
Let
\[
 \mathbf m_{\rm ext}=(\vpp pM)_{p\ne2,3},
 \qquad
 \mathbf a_{\rm ext}=(\vpp pA)_{p\ne2,3}.
\]
Only finitely many coordinates are nonzero.

\begin{proposition}[Kernel formula and dimension]\label{prop:K-dimension}
One has
\[
 K_x=\ker\left(
   \Q^2\longrightarrow\Q^{(\mathbb P\setminus\{2,3\})},
   \ (r,s)\longmapsto r\mathbf m_{\rm ext}+s\mathbf a_{\rm ext}
 \right).
                                                        \tag{8.2}\label{eq:K-valuations}
\]
Moreover:
\begin{enumerate}
\item $\dim_\Q K_x\le1$;
\item $K_x$ contains no vector $(r,s)$ with $rs>0$;
\item $K_x=\Q(1,0)$ exactly when $M$ is smooth;
\item $K_x=\Q(0,1)$ exactly when $A$ is smooth.
\end{enumerate}
\statusnew
\end{proposition}

\begin{proof}
Membership of $M^rA^s$ in $\Gamma$ is exactly the vanishing of every prime valuation outside
$2,3$, giving \eqref{eq:K-valuations}.  At least one of the two external vectors is nonzero by
\cref{prop:not-both-smooth}; hence the linear map has rank at least one and the kernel has
dimension at most one.

If $r,s>0$, any external prime occurring in $M$ or $A$ has positive valuation in $M^rA^s$; if
$r,s<0$, it has negative valuation.  Thus no same-sign vector lies in the kernel.  Finally, if
$M$ is smooth then the first external column is zero and the second is nonzero, so the kernel
is the first coordinate axis; the converse follows by testing $(1,0)$.  The statement for
$A$ is symmetric.
\end{proof}

\subsection{The kernel is independent of the nonintegral solution}

The exact conditional monoid already suggests that the same obstruction should persist under
translations and positive integer multiples.  The algebraic spectrum gives a stronger and
more conceptual result.

\begin{theorem}[Global kernel invariance]\label{thm:K-invariant}
Let $x,y\in\Sol\setminus\Z$.  Then
\[
  K_x=K_y.
\]
For $q=(r,s)\in\Q^2$, the following are equivalent:
\begin{enumerate}
\item $q\in K_x$;
\item $q\in K_y$;
\item $(2^r3^s)^{xy}$ is algebraic.
\end{enumerate}
If these conditions fail, then
\[
  xy\,(r\log2+s\log3)\notin\LL.
                                                        \tag{8.3}\label{eq:strong-mixed}
\]
\statusnew
\end{theorem}

\begin{proof}
Put $\alpha=2^r3^s\in\Gamma$.  Then $B=\alpha^x$ is algebraic.  By the exact spectrum applied
to the nonintegral solution $y$,
\[
  \alpha^{xy}=B^y\in\Qbar
  \quad\Longleftrightarrow\quad
  B\in\Gamma
  \quad\Longleftrightarrow\quad
  q\in K_x.
\]
The left side is symmetric in $x,y$, so the same condition is equivalent to $q\in K_y$.
If $B\notin\Gamma$, the strong spectrum applied to base $B$ and exponent $y$ gives
\[
 y\log B=xy\log\alpha\notin\LL.
\]
\end{proof}

Hence, under failure, there is a single global subspace $K$, not one kernel per solution.  It
measures exactly which rational directions in the $2,3$-radical group survive one additional
nonintegral exponentiation.

\begin{corollary}[Strong mixed-base extension]
Let $x,y$ be nonintegral solutions and let $r,s\in\Q$ have the same nonzero sign.  Then
\[
  (2^r3^s)^{xy}
\]
is log-strongly transcendental.  In particular $6^{xy}$ is log-strongly transcendental.
\end{corollary}

The Lean theorem
\lean{TwoBaseNonintegerSolution.transcendental_mixed_two_three_squared_exponents}
already gives ordinary transcendence for positive integers $r,s$ at $y=x$.  The corollary
above extends it to rational same-sign directions, to arbitrary pairs of nonintegral
solutions $x,y$, and to the log-strong conclusion \eqref{eq:strong-mixed}.

\subsection{M\"obius-orbit characterization}

\begin{theorem}[M\"obius criterion]\label{thm:mobius}
For $x\in\Sol\setminus\Z$, the following are equivalent:
\begin{enumerate}
\item $K_x\ne\{0\}$;
\item there exist $u,v,r,s\in\Q$ with
      \[
        us-vr\ne0,
        \qquad
        x=\frac{u+v\vartheta}{r+s\vartheta};             \tag{8.4}\label{eq:mobius}
      \]
\item $x\in\PGL_2(\Q)\cdot\vartheta$;
\item the four numbers $1,\vartheta,x,x\vartheta$ are linearly dependent over $\Q$.
\end{enumerate}
If these conditions hold, their $\Q$-span has dimension exactly three; otherwise it has
dimension four.  \statusnew
\end{theorem}

\begin{proof}
If $(r,s)\in K_x\setminus\{0\}$, write
\[
  M^rA^s=2^u3^v
\]
with $u,v\in\Q$.  Taking logarithms and dividing by $\log2$ gives
\[
  x(r+s\vartheta)=u+v\vartheta.
\]
The denominator cannot vanish because $\vartheta$ is irrational.  If
$us-vr=0$, the numerator and denominator are rationally proportional and $x$ is rational,
contrary to \cref{lem:rational-integral}.  This proves (1)$\Rightarrow$(2), and
(2)$\Rightarrow$(1) follows by exponentiating the same identity.

Equation \eqref{eq:mobius} is precisely membership in the rational projective-linear orbit.
After cross-multiplication it is a rational linear relation among
$1,\vartheta,x,x\vartheta$.  Conversely, any such relation must involve $x$ or
$x\vartheta$ because $1,\vartheta$ are independent, and rearrangement gives
\eqref{eq:mobius}.  Finally, two independent relations would make $K_x$ two-dimensional,
contradicting \cref{prop:K-dimension}; thus the dependent case has dimension exactly three.
\end{proof}

This is a genuinely different reformulation from the usual four-exponentials matrix.  It
asks whether the hypothetical exponent is a rational fractional-linear transform of the
fixed transcendental number $\vartheta$.  The answer is not forced by the original
hypotheses: the generic kernel branch has no such relation.

\subsection{Four counterexample types}\label{sec:kernel-types}

The external valuation vectors give an exhaustive four-way classification.

\begin{theorem}[Kernel-type classification]\label{thm:four-types}
Let $x\in\Sol\setminus\Z$.  Exactly one of the following occurs.

\begin{enumerate}[label=\textbf{Type \Roman*.}]
\item \textbf{$M$ smooth.}
      $M=2^a3^b$ with $b>0$, $K=\Q(1,0)$,
      $x=a+b\vartheta$, $E_3$ is algebraic, and $E_2$ is log-strongly transcendental.

\item \textbf{$A$ smooth.}
      $A=2^c3^d$ with $c>0$, $K=\Q(0,1)$,
      $x=d+c\eta$, $E_2$ is algebraic, and $E_3$ is log-strongly transcendental.

\item \textbf{Common external valuation ray.}
      Both outputs are nonsmooth, but their external valuation vectors are proportional.
      There are coprime $d,e\in\N$ and a nonzero integer vector $\mathbf c\ge0$ such that
      \[
        \mathbf m_{\rm ext}=d\mathbf c,
        \qquad
        \mathbf a_{\rm ext}=e\mathbf c.
      \]
      Then
      \[
        K=\Q(e,-d),
      \]
      both $E_3,E_2$ are log-strongly transcendental, and for some $u,v\in\Z$,
      \[
        \frac{M^e}{A^d}=2^u3^v,
        \qquad
        x=\frac{u+v\vartheta}{e-d\vartheta}.             \tag{8.5}\label{eq:common-core-mobius}
      \]

\item \textbf{Independent external directions.}
      The two external valuation vectors are linearly independent.  Then $K=\{0\}$,
      $x\notin\PGL_2(\Q)\cdot\vartheta$, the four numbers
      $1,\vartheta,x,x\vartheta$ are $\Q$-linearly independent, and both
      $E_3,E_2$ are log-strongly transcendental.
\end{enumerate}
\statusnew
\end{theorem}

\begin{proof}
If one external vector vanishes, \cref{prop:K-dimension,thm:E3,thm:E2} give Types I and II;
they are mutually exclusive.  If both are nonzero, they have rank one or two.  In rank one,
clear denominators and divide the two scale factors by their gcd to obtain the displayed
coprime $d,e$ and a common nonnegative integer vector $\mathbf c$.  The kernel is then
$\Q(e,-d)$.  The external valuations cancel, so
$M^e/A^d\in\Gamma\cap\Q_{>0}$.  Its $2$- and $3$-valuations are integers,
giving $u,v\in\Z$.  Taking logarithms gives \eqref{eq:common-core-mobius}.  Rank two gives
zero kernel.  The claims about $E_3,E_2$ follow because both outputs are nonsmooth.
\end{proof}

Type III says that outside the distinguished primes, $M$ and $A$ are powers of one common
integer core:
\[
  M=2^{a_2}3^{a_3}C^d,
  \qquad
  A=2^{b_2}3^{b_3}C^e
\]
with $a_2,a_3,b_2,b_3\in\Nzero$ and $C$ supported only on primes outside $2,3$.  Type IV,
by contrast, supplies two external primes $p,q$ with nonzero valuation determinant
\[
 \Delta_{p,q}=
 \vpp pM\vpp qA-\vpp qM\vpp pA\ne0.                  \tag{8.6}\label{eq:external-minor}
\]
This distinction is central to the proposed local determinant program.

\section{Depth-three tower rigidity}\label{sec:tower}

The kernel gives a complete answer to the following question: starting from a positive
algebraic base $\alpha$, how many times can one repeatedly apply the same nonintegral
exponent $x$ before transcendence is forced?

For positive bases there is no ambiguity between iterated powers and powers of powers:
\[
  \bigl((\alpha^x)^x\bigr)^x
  =\alpha^{x^3}.
\]

\subsection{No nontrivial two-step return inside \texorpdfstring{$\Gamma$}{Gamma}}

For $\alpha=2^r3^s\in\Gamma$ define $T_x(\alpha)=\alpha^x$.  The map takes $\Gamma$ into
$\Apos$, and $K_x$ is exactly the set of elements whose image returns to $\Gamma$.

\begin{lemma}[No two-step internal orbit]\label{lem:no-two-step}
Let $x\in\Sol\setminus\Z$ and let $\alpha\in\Gamma\setminus\{1\}$.  If
$\alpha\in K_x$, then
\[
  T_x(\alpha)\notin K_x.
\]
Equivalently,
\[
 \alpha^x\in\Gamma
 \quad\Longrightarrow\quad
 \alpha^{x^2}\notin\Gamma.
\]
\statusnew
\end{lemma}

\begin{proof}
Suppose both $\alpha$ and $\beta=T_x(\alpha)$ lie in the one-dimensional space
$K_x\setminus\{0\}$.  Under the identification $\Gamma\cong\Q^2$, two nonzero elements on the
same rational line differ by a rational scalar, so $\beta=\alpha^q$ for some $q\in\Q$.
But $\beta=\alpha^x$ and $\alpha\ne1$, hence $x=q$, contradicting the irrationality of $x$.
Thus $\beta\notin K_x$, which means $T_x(\beta)=\alpha^{x^2}\notin\Gamma$.
\end{proof}

The lemma has a useful dynamical interpretation.  A nonzero direction in $K$ may survive one
return to the algebraic spectrum, but the returned direction cannot itself survive.  There
are no nontrivial cycles, fixed rays, or length-two paths entirely inside $\Gamma$.

\subsection{Exact survival depth}

\begin{theorem}[Depth-three algebraic-tower theorem]\label{thm:depth-three}
Let $x\in\Sol\setminus\Z$ and let $\alpha\in\Apos\setminus\{1\}$.  At least one of
\[
  \alpha^x,\qquad \alpha^{x^2},\qquad \alpha^{x^3}
\]
is transcendental.  More precisely, the first transcendental term occurs at exactly one of
the following depths:

\begin{enumerate}
\item \textbf{Depth 1:} if $\alpha\notin\Gamma$, then $\alpha^x$ is log-strongly
      transcendental.
\item \textbf{Depth 2:} if $\alpha\in\Gamma$ but $\alpha\notin K_x$, then
      $\alpha^x\in\Qbar\setminus\Gamma$ and $\alpha^{x^2}$ is log-strongly transcendental.
\item \textbf{Depth 3:} if $\alpha\in K_x\setminus\{1\}$, then
      $\alpha^x\in\Gamma$, $\alpha^{x^2}\in\Qbar\setminus\Gamma$, and
      $\alpha^{x^3}$ is log-strongly transcendental.
\end{enumerate}

Equivalently, for some $j\in\{1,2,3\}$,
\[
  x^j\log\alpha\notin\LL.
\]
\statusnew
\end{theorem}

\begin{proof}
Depth 1 is \cref{thm:strong-spectrum}.  Suppose $\alpha\in\Gamma$.  Then
$\beta=\alpha^x$ is algebraic.  If $\alpha\notin K_x$, then $\beta\notin\Gamma$; applying the
strong spectrum to the algebraic base $\beta$ and exponent $x$ gives
\[
 \alpha^{x^2}=\beta^x
\]
log-strongly transcendental.

If $\alpha\in K_x\setminus\{1\}$, then $\beta\in\Gamma$.  Thus
$\gamma=\beta^x=\alpha^{x^2}$ is algebraic.  By \cref{lem:no-two-step},
$\beta\notin K_x$, hence $\gamma\notin\Gamma$.  A final application of the strong spectrum to
base $\gamma$ gives $\gamma^x=\alpha^{x^3}$ log-strongly transcendental.
\end{proof}

This theorem is a compact synthesis of the auxiliary-base spectrum and the kernel.  It also
places a strict ceiling on how far an algebraic power tower can imitate the integral first
level of a counterexample.

\subsection{Exact behavior of bases \texorpdfstring{$2$, $3$, and $6$}{2, 3, and 6}}

\begin{corollary}[Base $2$ tower dichotomy]\label{cor:base2-tower}
For a nonintegral solution $x$ with $M=2^x$:
\begin{enumerate}
\item if $M$ is not $\{2,3\}$-smooth, then $2^{x^2}$ is log-strongly transcendental;
\item if $M$ is $\{2,3\}$-smooth, then $2^{x^2}$ is algebraic (indeed an integer) and
      $2^{x^3}$ is log-strongly transcendental.
\end{enumerate}
\end{corollary}

\begin{proof}
The vector $(1,0)$ lies in $K_x$ exactly when $M\in\Gamma$, equivalently when the integer $M$
is smooth.  Apply \cref{thm:depth-three}.  In the smooth case, write
$M=2^a3^b$; then
\[
  2^{x^2}=M^x=(2^x)^a(3^x)^b=M^aA^b\in\N.
\]
\end{proof}

The symmetric statement holds for base $3$ and $A$.

\begin{corollary}[Why base $6$ always fails at depth two]
For every nonintegral solution $x$,
\[
  x^2\log6\notin\LL,
\]
so $6^{x^2}$ is log-strongly transcendental.
\end{corollary}

\begin{proof}
The coefficient vector $(1,1)$ has positive coordinate product and therefore does not lie in
$K_x$.  Since $6\in\Gamma$, it has exact survival depth two.
\end{proof}

This recovers the current base-$6$ theorem and places it inside the kernel-verified family
that every $(2^i3^j)^{x^2}$ with $i,j>0$ is transcendental.  The new strengthening is the
log-strong conclusion, its extension to arbitrary rational same-sign directions and pairs of
nonintegral solutions, and the exact explanation of why the product base avoids the
exceptional smooth axes.

\subsection{New one-base tower equivalences of Alaoglu--Erd\H{o}s}

\begin{theorem}[Universal three-level detector]\label{thm:universal-tower-detector}
Fix any $\alpha\in\Apos\setminus\{1\}$.  The Alaoglu--Erd\H{o}s conjecture is equivalent to
\[
 \forall x\in\R,\quad
 2^x,3^x\in\Z
 \Longrightarrow
 \alpha^x,\alpha^{x^2},\alpha^{x^3}\in\Qbar.
                                                        \tag{9.1}\label{eq:three-level-detector}
\]
The same equivalence holds with membership in $\LL$ required of the three logarithms
$x^j\log\alpha$.  \statusnew
\end{theorem}

\begin{proof}
If the conjecture holds, every solution is a nonnegative integer, and all three powers are
algebraic.  If it fails, apply \cref{thm:depth-three} to a nonintegral solution.
\end{proof}

For bases in $\Gamma$, the first level is automatically algebraic.  This gives a particularly
simple new reformulation.

\begin{corollary}[Base-$2$ cubic-tower equivalence]\label{cor:base2-equivalence}
The Alaoglu--Erd\H{o}s conjecture is equivalent to
\[
 \forall x\in\R,\quad
 2^x,3^x\in\Z
 \Longrightarrow
 2^{x^2}\in\Qbar\ \text{and}\ 2^{x^3}\in\Qbar.
                                                        \tag{9.2}\label{eq:base2-cubic}
\]
The same statement holds with base $3$ in place of base $2$.
\end{corollary}

This differs from the repository's existing equivalent formulation requiring both
$2^{x^2}$ and $3^{x^2}$ to be algebraic.  A single original base suffices if one allows one
more tower level.

\subsection{A compact survival-depth table}

For a nonintegral solution, the first forced transcendence depth of a positive algebraic
base $\alpha\ne1$ is:

\begin{center}
\begin{tabular}{lll}
\toprule
Location of $\alpha$ & First log outside $\LL$ & Structural reason \\
\midrule
$\alpha\notin\Gamma$ & $x\log\alpha$ & six-exponential spectrum \\
$\alpha\in\Gamma\setminus K$ & $x^2\log\alpha$ & first image exits $\Gamma$ \\
$\alpha\in K\setminus\{1\}$ & $x^3\log\alpha$ & no two-step internal orbit \\
\bottomrule
\end{tabular}
\end{center}

This table is exact, not merely an upper bound.

\section{A catalogue of equivalent formulations and near-equivalences}

The corpus and the new results now provide a fairly rich dictionary.  Let \textsf{AE} denote
the original conjecture.

\begin{longtable}{@{}>{\raggedright\arraybackslash}p{0.25\textwidth}>{\raggedright\arraybackslash}p{0.48\textwidth}>{\raggedright\arraybackslash}p{0.15\textwidth}@{}}
\toprule
Formulation & Statement & Status \\
\midrule
\endhead
Original & $2^x,3^x\in\Z\Rightarrow x\in\Z$. & Open \\
Candidate form & $M^{\vartheta}\in\N\Rightarrow M$ is a power of $2$. & Exact \\
Localized radical form & No odd power-primitive $u>1$ has
$u^{\vartheta}\in\Z[1/3]$ with the canonical side conditions. & Lean/paper \\
No nontrivial generator & The conditional least generator $\beta$ does not exist. & Lean \\
Base-$5$ algebraicity & Every solution has $5^x\in\Qbar$. & Lean \\
Full algebraic output locus & For every positive algebraic $\gamma$,
$\gamma^x\in\Qbar$ iff $x\in\Z$ or $\gamma\in\Gamma$. & Lean \\
Arbitrary algebraic detector & For any fixed $\alpha\in\Apos\setminus\Gamma$, every solution
has $\alpha^x\in\Qbar$. & New corollary of Lean locus \\
Both first iterates & Every solution has both $2^{x^2},3^{x^2}\in\Qbar$. & Lean \\
Base-$6$ first iterate & Every solution has $6^{x^2}\in\Qbar$. & Lean \\
Ratio-squared rank two & The algebraic locus among $(2^i/3^j)^{x^2}$ contains two
noncollinear exponent vectors. & Lean \\
Single-base cubic tower & Every solution has $2^{x^2},2^{x^3}\in\Qbar$. & New \\
Universal algebraic tower & For one fixed $\alpha\ne1$, all three of
$\alpha^x,\alpha^{x^2},\alpha^{x^3}$ are algebraic for every solution. & New \\
Strong detector & For fixed $\alpha\notin\Gamma$, every solution has
$x\log\alpha\in\LL$. & New \\
Subquadratic cumulative count & The candidate count is $o((\log X)^2)$ in the appropriate
logarithmic scale, excluding a rank-two conditional monoid. & Paper equivalence \\
Sublinear block count & A sufficiently strong $o(\log X)$ dyadic candidate bound, with
uniformity adequate to exclude the exact monoid. & Near-equivalence \\
Nonlinear closure & Closure of $\Sol$ under multiplication or squaring would imply AE, since
nonintegral products are excluded under failure. & Paper implication \\
Four exponentials & The four exponentials conjecture applied to $(1,x)$ and
$(\log2,\log3)$. & Classical implication \\
\bottomrule
\end{longtable}

A few distinctions matter:

\begin{itemize}
\item The detector equivalences are truly pointwise: one fixed independent algebraic base is
enough.
\item The tower equivalence works even for bases inside $\Gamma$, where first-level
algebraicity is automatic under a counterexample.
\item A generic $O(\log X)$ block bound is not automatically contradictory; its leading
constant and uniformity must beat the conditional monoid.  An $o(\log X)$ bound is cleanly
sufficient.
\item Closure under a nonlinear operation is not known; the conditional theorem says exactly
why obtaining it would be decisive.
\end{itemize}

\section{Finite verification and computational audit}\label{sec:computation}

\subsection{Formal and nonformal finite frontiers}

Two finite statements must be kept distinct.

\begin{enumerate}
\item \statuslean{} The current formal development proves that no nonintegral solution occurs
      with $2^x<4096=2^{12}$.  Thus a counterexample has $x\ge12$.
\item \statuscomp{} The third article~\cite{ProveItReport3} reports an exhaustive directed-rounding scan of every
      positive integer $M<2^{26}$ and finds no candidate other than powers of $2$.  This gives
      the computational conclusion $x>26$ for a counterexample.
\end{enumerate}

The second result is stronger numerically but weaker epistemically.  It should not be cited as
a theorem of the Lean kernel.

\subsection{What the MPFR scan must certify}

For each $M$, one must decide whether
\[
  M^{\vartheta}=\exp\!\left(\frac{\log M\,\log3}{\log2}\right)
\]
is an integer.  Ordinary floating-point evaluation followed by a tolerance test is
insufficient: a missed candidate could lie closer to an integer than the tolerance, and
rounding error could cross the integer boundary.

The reported program instead constructs a rigorous interval enclosure using directed MPFR
rounding.  Conceptually, it computes lower and upper bounds for
\[
  y=\vartheta\log M,
  \qquad
  z=e^y,
\]
and verifies that the resulting interval contains no integer.  A fast approximate locator
selects nearby integer candidates, but the final accept/reject test uses the outward-rounded
interval, not the approximate value.  The code widens the locator neighborhood to protect
against a one-unit mislocation.

\subsection{Source-level audit findings}

The source logic described in the third report has the right shape:

\begin{itemize}
\item lower and upper logarithms use opposite rounding directions;
\item multiplication by the positive ratio $\log3/\log2$ preserves endpoint order;
\item exponentiation is monotone and also outward rounded;
\item every integer that could intersect the final interval is checked;
\item exact powers of $2$ are recognized as expected candidates rather than false failures;
\item the loop bounds cover $1\le M<2^{26}$.
\end{itemize}

I found no conceptual rounding error in that design.  However, a source audit is not an
execution transcript, and the report's binary/hash record is not a compact independently
checkable certificate of all $2^{26}$ interval exclusions.

\subsection{How to turn the scan into a stronger artifact}

A preferable next finite artifact would split the range into blocks and emit, for each block,
a compact certificate containing:

\begin{itemize}
\item the MPFR precision and exact constants used;
\item a hash of the source and compiler command;
\item the minimum certified distance from an interval to the nearest integer;
\item the input attaining that minimum;
\item an independently replayable block checksum;
\item optionally, rational interval endpoints for only the extremal near misses.
\end{itemize}

A checker substantially smaller than the search program could then replay the certificates.
The ultimate formal route is a Lean interval proof for blocks using rational bounds for
$\log$ and $\exp$, but extending the theoretical structure is currently more valuable than
pushing the raw cutoff by another few bits.

\section{Lean formalization roadmap}\label{sec:lean-roadmap}

The repository advanced during this audit.  The current snapshot has already formalized the
ordinary algebraic-output layer that an earlier draft of this report had listed as future
work.  The remaining roadmap therefore separates (i) newly available infrastructure,
(ii) elementary algebraic consequences that can be added now, and (iii) genuinely stronger
logarithmic transcendence statements requiring new classical inputs.

\subsection{What the current snapshot just formalized}

The new kernel-verified layer consists of:

\begin{itemize}
\item \lean{AlgebraicSixExponentials}: a symmetric positive-real algebraic
      six-exponentials interface;
\item \lean{AlgebraicRationalBase} and \lean{AlgebraicBaseUnitPower}: exact rational and
      arbitrary positive-algebraic output loci;
\item \lean{AlgebraicRankThreeEquivalence}: the rank-three algebraic-output reformulation
      of Alaoglu--Erd\H{o}s;
\item \lean{AlgebraicOutputLocus}: mixed squared-output transcendence, the rank-one ratio
      locus, and an explicit external-prime valuation line;
\item \lean{CanonicalRadicalAlgebraicOutput}, \lean{CanonicalRadicalDivisibleHull}, and
      \lean{CanonicalRadicalFieldNorm}: corresponding upgrades for the canonical radicals.
\end{itemize}

Consequently the ordinary $2,3$-saturated algebraic spectrum is no longer a formalization
proposal.  The genuinely new formal targets in this report are the Smith refinement, the
full rational kernel and its M\"obius classification, the depth-three orbit theorem, and the
log-strong upgrades.

\subsection{What can plausibly be formalized now}

\subsubsection{A base-\texorpdfstring{$2$}{2} squared/cubed tower theorem}

The base-$2$ specialization of \cref{cor:base2-tower} only needs the current natural-base
algebraicity classification, the smooth-output decomposition, and transcendence of
$\vartheta$.  It does not require arbitrary algebraic bases or Roy's theorem.  A suitable new
module could be
\[
 \texttt{TwoBaseIntegerExponent/SquaredCubedTower.lean}.
\]
The target interface is schematically:

\begin{lstlisting}[style=leanstyle]
/-- A nonintegral two-base solution cannot have algebraic
    base-2 outputs at both the squared and cubed exponents. -/
theorem TwoBaseNonintegerSolution.not_both_two_squared_cubed_algebraic
    {x : R} (hx : TwoBaseNonintegerSolution x) :
    not (IsAlgebraic Q ((2 : R) ^ (x * x)) and
         IsAlgebraic Q ((2 : R) ^ (x * x * x))) := by
  ...

/-- A one-base, two-level reformulation of Alaoglu--Erdos. -/
theorem alaogluErdosConjecture_iff_two_squared_cubed_algebraic :
    AlaogluErdosConjecture <->
      forall x : R,
        (exists z : Z, (z : R) = (2 : R) ^ x) ->
        (exists z : Z, (z : R) = (3 : R) ^ x) ->
        IsAlgebraic Q ((2 : R) ^ (x * x)) and
        IsAlgebraic Q ((2 : R) ^ (x * x * x)) := by
  ...
\end{lstlisting}

A direct proof splits on whether $M=2^x$ is $\{2,3\}$-smooth.  In the nonsmooth case, the
current natural-base theorem makes $M^x=2^{x^2}$ transcendental.  In the smooth case,
$2^{x^2}$ is an integer.  If it were also smooth, expanding
$x=a+b\vartheta$ and $(a+b\vartheta)^2$ would give a nonzero quadratic polynomial for
$\vartheta$; hence the integer base $2^{x^2}$ is nonsmooth, and its $x$-th power
$2^{x^3}$ is transcendental.

\subsubsection{External valuation rank and the four types}

A module
\[
 \texttt{TwoBaseIntegerExponent/ExternalValuationKernel.lean}
\]
should build directly on \lean{AlgebraicOutputLocus}, which already extracts one external
prime and proves a rank-one nonnegative-lattice slice.  The new module can avoid quotient
vector spaces initially.  For natural outputs $M,A$, define the predicate
that all external valuations of $M^rA^s$ cancel after clearing a common denominator.  The
rank-zero/one/two split can be expressed through proportionality of finitely supported
valuation functions and the existence of a nonzero $2\times2$ minor.  Useful target theorems
are:

\begin{lstlisting}[style=leanstyle]
/-- Both nonsmooth outputs either share one external valuation ray
    or have a nonzero external two-prime minor. -/
theorem common_external_ray_or_exists_nonzero_minor ... :
  (exists d e C, ...) or
  (exists p q : Nat, p.Prime and q.Prime and
    p != 2 and p != 3 and q != 2 and q != 3 and
    padicValNat p M * padicValNat q A !=
      padicValNat q M * padicValNat p A) := by
  ...
\end{lstlisting}

Mathlib's finitely supported functions and unique factorization API should suffice.  The main
engineering issue is choosing a valuation representation that handles zero entries without
repeated coercion friction.

\subsubsection{Nonlinear closure from the third report}

The polynomial identity argument is particularly formalization-friendly.  Under failure,
write every solution uniquely as $n+k\beta$.  If
\[
 (n+k\beta)(m+\ell\beta)=r+s\beta,
\]
then transcendence of $\beta$ forces $k\ell=0$.  A module
\[
 \texttt{TwoBaseIntegerExponent/NonlinearClosure.lean}
\]
could include:

\begin{itemize}
\item product closure iff one factor is integral;
\item $P(\beta)\in\Sol$ iff $P$ is affine with nonnegative integral coefficients at $\beta$;
\item the rational-function field intersection;
\item exact common-perfect-power degree $\gcd(n,k)$;
\item primitive-pair counting after importing the necessary M\"obius-sum estimates.
\end{itemize}

The first three need little analytic number theory and would materially tighten the formal
conditional picture.

\subsubsection{Smith-normal-form rational-output lattices}

The valuation criterion \cref{thm:rational-integral-spectrum} is elementary once the
algebraic base has already been represented as $2^r3^s$.  A first formal milestone can state
it for rational pairs with an explicit common denominator, avoiding a general positive
algebraic-number API.  Mathlib already contains Smith normal form infrastructure in various
matrix/module guises, but a bespoke rank-two theorem may be shorter and more stable.

\subsection{What needs a larger transcendence formalization}

\subsubsection{General saturated-pair packaging}

The specialized $2,3$ statement is now kernel-verified.  A lightweight optional module can
package the same idea for arbitrary multiplicatively independent positive algebraic bases
$a,b$:
\[
 \alpha^x\in\Qbar\quad\Longleftrightarrow\quad
 \alpha\in\GammaSat(a,b),
\]
under $x\notin\Q$ and algebraicity of $a^x,b^x$.  The existing generic theorem in
\lean{AlgebraicSixExponentials} supplies the difficult implication; the reverse implication
is elementary rational-power algebra.  This is primarily an API/generalization task, not a
new determinant project.

\subsubsection{The strong spectrum and symmetric constants}

The assertions involving $\LL$ require formal versions of:

\begin{itemize}
\item Baker's homogeneous linear independence theorem for logarithms of algebraic numbers;
\item Roy's rank theorem/strong six exponentials theorem;
\item a coherent formal definition of the $\Qbar$-linear logarithm span, including branches.
\end{itemize}

This is a substantial project in its own right.  A pragmatic layering is to formalize the
ordinary-transcendence versions first; every kernel and tower statement has a non-strong
version using only the exact algebraic spectrum.

\subsection{Suggested module order}

A dependency-minimizing order is:

\begin{enumerate}
\item \texttt{NonlinearClosure.lean};
\item \texttt{SquaredCubedTower.lean};
\item \texttt{ExternalValuationKernel.lean};
\item \texttt{AlgebraicRadicalOutputLattice.lean};
\item \texttt{GeneralSaturatedBaseSpectrum.lean} (optional API packaging);
\item \texttt{StrongAlgebraicBaseSpectrum.lean};
\item \texttt{SymmetricTowerConstants.lean}.
\end{enumerate}

The first five reuse the current determinant and algebraic-output infrastructure.  Only the
last two require a formal Baker/Roy logarithm-span layer.

\subsection{Proof-audit checklist for new formal modules}

Each proposed theorem should explicitly track:

\begin{itemize}
\item positivity hypotheses required by real-power multiplication laws;
\item whether ``integer'' is represented by $\Z$ or positive $\N$;
\item the distinction between rational powers and natural powers;
\item branch-free real logarithms versus complex logarithm spans;
\item finite prime support before applying gcd or Smith arguments;
\item the exact source of irrationality/transcendence of $x$ and $\vartheta$;
\item status boundaries between kernel proof, imported classical theorem, and interval
      computation.
\end{itemize}

\section{Branch-sensitive future research program}\label{sec:future}

A single undifferentiated instruction to ``use more transcendence theory'' is not a research
plan.  The kernel classification suggests four different hypothetical worlds, with different
arithmetic resources.  The most promising next work should attack them separately and then
recombine the results.

\subsection{Priority 1: a two-prime tropical/\texorpdfstring{$p$-adic}{p-adic} determinant for Type IV}

Type IV supplies external primes $p,q$ with the nonzero valuation minor
\eqref{eq:external-minor}.  This is information discarded by the current archimedean
interpolation determinants.

For conditional lattice points $(n_i,k_i)$ define synchronized integer outputs
\[
  m_i=2^{n_i}M^{k_i},
  \qquad
  a_i=3^{n_i}A^{k_i}.
\]
Choose column exponent pairs $(u_j,v_j)\in\Nzero^2$ and form
\[
  D=\det\bigl(m_i^{u_j}a_i^{v_j}\bigr)_{0\le i,j<r}.
                                                        \tag{13.1}\label{eq:adelic-det}
\]
At an external prime $p$,
\[
  \vpp p{m_i^{u_j}a_i^{v_j}}
   =k_i\bigl(u_j\vpp pM+v_j\vpp pA\bigr).
                                                        \tag{13.2}\label{eq:tropical-weight}
\]
Thus every determinant term has a tropical assignment weight.  Ordering the $k_i$ and the
column slopes
\[
 c_{p,j}=u_j\vpp pM+v_j\vpp pA
\]
allows the rearrangement inequality to identify a minimal valuation permutation.  A unique
minimum, together with a nonzero residual determinant modulo $p$, gives an exact or sharp
lower bound for $v_p(D)$.  The independent $q$-slope ordering supplied by
$\Delta_{p,q}\ne0$ should provide a second nonarchimedean gain that cannot be absorbed into
the first.

The target is an adelic product-formula contradiction:
\[
  |D|_\infty
  \prod_{\ell\in\{p,q\}}|D|_\ell
  <1.
\]

after all remaining norms are bounded by $1$, while $D$ is a nonzero integer.  The
archimedean factor comes from clustered nodes; the $p$- and $q$-adic factors come from the
exact $(n,k)$ lattice and external valuation directions.

\begin{problem}[Tropical noncancellation lemma]
Choose a family of exponent pairs $(u_j,v_j)$ for which the minimal assignment in
\eqref{eq:tropical-weight} is unique at $p$ and at $q$, or prove that the leading residual
minor is a unit after a controlled column transformation.
\end{problem}

\begin{problem}[One-logarithm recovery]
Show that the combined nonarchimedean gain improves the unconditional dyadic candidate count
from $O((\log X)^2)$ to $o(\log X)$ for configurations having a nonzero external valuation
minor.
\end{problem}

This is the highest-priority direction because Type IV is the generic branch and because it
uses genuinely new arithmetic information rather than only enlarging an archimedean
exponent box.

\subsection{Priority 2: isolate and attack the common-core Type III branch}

In Type III one has, after normalization,
\[
  M=2^{a_2}3^{a_3}C^d,
  \qquad
  A=2^{b_2}3^{b_3}C^e,
\]
with one external primitive core $C>1$ and coprime $d,e>0$.  The relation
\[
  M^e/A^d=2^u3^v
\]
produces the M\"obius identity \eqref{eq:common-core-mobius}.

This branch has no two-prime valuation area: every external valuation minor vanishes.
Instead, it has an exact common-core quotient and a rational linear relation among
$1,\vartheta,x,x\vartheta$.

Promising targets are:

\begin{enumerate}
\item derive a canonical normalization of $C,d,e,a_i,b_i$ compatible with the repository's
      primitive output cores and prove uniqueness;
\item translate all perfect-power and distinguished-base depth inequalities into constraints
      on $d,e,u,v$ in \eqref{eq:common-core-mobius};
\item study whether the denominator $e-d\vartheta$ can be forced to give too good a rational
      approximation to $\vartheta$ once $M,A$ are integers of the prescribed sizes;
\item build a one-core interpolation determinant over $\Q(C)$ whose conjugate or local norms
      retain information lost when all external vectors are proportional;
\item seek a reduction from Type III to one of the smooth axis types by proving that a
      power-primitive common core cannot occur on both output sides.
\end{enumerate}

\begin{problem}[Common-core exclusion]
Prove that a nonintegral two-base solution cannot have proportional nonzero external valuation
vectors.  Equivalently, eliminate Type III or force one output to be $\{2,3\}$-smooth.
\end{problem}

Even a partial result such as excluding $d=e$, squarefree $C$, or a single external prime
would materially narrow the conjecture.

\subsection{Priority 3: pure-radical attacks on Types I and II}

In Type I,
\[
  x=a+b\vartheta,
  \qquad
  E_3^b=A/3^a\in\Q_{>0},
\]
so $E_3=3^{\vartheta}$ is an algebraic radical with explicitly bounded degree and known
valuation data.  Type II is symmetric for $E_2=2^{1/\vartheta}$.

Unlike the original relation, this is a legitimate number-field object.  One can place $E_3$
in a pure extension, analyze all conjugates, and use norms.  The current
\texttt{AlgebraicIntegerDeterminant} and \texttt{AlgebraicOutputBridge} modules are directly
relevant.

Concrete objectives are:

\begin{itemize}
\item prove that the exact degree of $E_3$ implied by the primitive normalization is
      incompatible with the $3$-adic valuation of $A/3^a$;
\item combine $2^{\vartheta}=3$ with the minimal polynomial of $E_3$ in a two-variable
      interpolation determinant using both real embeddings and finite places;
\item exploit the non-smoothness of $A$ to obtain an external prime in the radical's norm;
\item derive an algebraic auxiliary output outside $\Gamma$ from a conjugate or norm of
      $E_3$; the spectrum theorem would then close the argument immediately.
\end{itemize}

The bare statement that $E_3$ should be transcendental is itself a difficult quadratic
logarithm problem.  The useful extra input is not merely algebraicity but the exact radical
identity with an integer output and the canonical perfect-power constraints.

\begin{problem}[Axis radical contradiction]
Assume $M=2^a3^b$ with $b>0$ and $A=3^x\in\N$.  Use the canonical primitive normalization to
exclude the radical equality $(3^{\vartheta})^b=A/3^a$.
\end{problem}

\subsection{Priority 4: an exact-lattice interpolation determinant}

The current unconditional determinant estimates treat arbitrary clusters of candidates.  A
counterexample gives much more: every candidate has coordinates
\[
  m(n,k)=2^nM^k,
  \qquad n,k\in\Nzero.
\]
The dyadic constraint is a thin strip
\[
  \log X\le n\log2+k\log M\le\log(2X)
\]
through the integer lattice.  Consecutive points have an ordered Sturmian geometry controlled
by the continued fraction of $\beta=\log_2M$.

Rather than selecting arbitrary nodes, construct the determinant directly from a complete
strip segment and choose columns dual to its two lattice directions.  Desired gains include:

\begin{itemize}
\item exact products of row gaps from the continued-fraction ordering;
\item cancellation of the translation direction $n$ before estimating the determinant;
\item a determinant dimension proportional to the actual $O(\log X)$ strip count, not the
      square of an exponent-box side length;
\item simultaneous use of output coordinates $3^nA^k$ and external prime valuations.
\end{itemize}

\begin{problem}[Strip determinant]
Construct a nonzero integer determinant attached to all lattice points in one dyadic strip
whose archimedean upper bound has an exponent linear, rather than quadratic, in the strip
length.
\end{problem}

This route may merge naturally with the Type-IV adelic determinant.

\subsection{Priority 5: force an algebraic bridge outside \texorpdfstring{$\Gamma$}{Gamma}}

The exact spectrum theorem turns the whole conjecture into a bridge-construction problem:
under the original integer hypotheses, produce any positive algebraic
$\alpha\notin\Gamma$ with algebraic $\alpha^x$.

Potential bridge mechanisms include:

\begin{itemize}
\item a norm from an algebraic radical appearing in the canonical output normalization;
\item an eigenvalue or determinant of a matrix built functorially from $M,A$ whose $x$-th
      power is another algebraic matrix invariant;
\item a resultant linking the two integer recurrences $M^n$ and $A^n$;
\item an algebraic limit forced by a sufficiently strong synchronized rational approximation;
\item a local-to-global identity from two independent external primes.
\end{itemize}

The spectrum theorem also prevents wasted effort: a candidate bridge that simplifies to
$2^r3^s$ cannot solve the problem at first level.

\subsection{Priority 6: denominator-sensitive orbit methods}

Reducing $kx$ modulo $1$ gives the dense orbit
\[
 \left(2^{\{kx\}},3^{\{kx\}}\right)
 =\left(\frac{M^k}{2^{\lfloor kx\rfloor}},
        \frac{A^k}{3^{\lfloor kx\rfloor}}\right).
\]
The first coordinate is dyadic rational and the second triadic rational with synchronized
exponent denominators.  Generic rational-point counting on a transcendental curve is too
weak: $O(\log H)$ points up to height $H$ are compatible with known bounds.

A useful theorem must exploit the restricted denominator supports and their shared orbit
index.  Possible frameworks are:

\begin{itemize}
\item a quantitative $S$-integral Pila--Wilkie theorem with denominator support;
\item a Subspace-Theorem estimate for local Taylor approximants to the curve
      $v=u^{\vartheta}$;
\item an adelic determinant on the orbit points, retaining $2$- and $3$-adic denominators;
\item a discrepancy estimate weighted by exact denominator growth rather than only orbit
      count.
\end{itemize}

\subsection{Priority 7: computation as a certificate generator, not a substitute}

Extending the scan from $2^{26}$ to $2^{30}$ or $2^{32}$ is straightforward engineering but
unlikely to alter the theoretical bottleneck.  More valuable computational work would:

\begin{itemize}
\item locate extremal near-candidates and test conjectured lower-bound shapes;
\item search for the best determinant column sets by integer optimization;
\item explore tropical valuation orderings for Type IV;
\item enumerate possible small canonical normalizations $(a,b,d,e)$ consistent with all
      formal inequalities;
\item produce replayable interval certificates.
\end{itemize}

\subsection{Recommended order of attack}

The practical order I recommend is:

\begin{enumerate}
\item formalize the base-$2$ squared/cubed theorem and nonlinear closure package;
\item formalize the four kernel types at the natural-output valuation level;
\item prototype the Type-IV two-prime tropical determinant computationally;
\item isolate the exact Type-III common-core normal form and prove special-case exclusions;
\item revisit Types I/II with the existing algebraic-integer determinant machinery;
\item only then invest in a larger generic determinant or a deeper finite scan.
\end{enumerate}

The branch split prevents a method suited to two independent external primes from being
judged against a one-core exceptional case where it cannot possibly work.

\section{Conclusion}

The conjecture remains unproved in this analysis.  The exact obstruction is visible and
honest: the original four logarithms form the unresolved singular $2\times2$ configuration
of four exponentials.  The ProveIt corpus has nevertheless converted a hypothetical failure
into an extraordinarily rigid object: a least transcendental generator, a free rank-two
solution/output monoid, exact counting laws, radical normal forms, strong spacing bounds, and
severe restrictions on all auxiliary outputs.

The main conceptual advance of the present report is that the auxiliary-output picture is now
complete for \emph{all positive algebraic bases}.  Under a counterexample, algebraic output is
possible exactly on the saturated $2,3$-radical group $\Gamma$.  Inside that group, prime
valuations classify rational and integral output; a global kernel of dimension at most one
classifies algebraic second iterates; and no nontrivial direction survives twice.  Hence every
algebraic-base tower becomes transcendental by the third level, yielding new one-base
reformulations of Alaoglu--Erd\H{o}s.

The kernel also divides the hypothetical world into four arithmetically distinct branches.
That is more than a taxonomy: it identifies the data each future method must use.  The generic
branch supplies a nonzero two-prime valuation area and invites an adelic tropical determinant.
The common-core branch supplies a rational M\"obius relation and a single primitive external
core.  The two axis branches turn one of the symmetric constants $3^{\log_2 3}$ and
$2^{\log_3 2}$ into an explicit algebraic radical.  Eliminating these branches separately is,
in my assessment, the most promising path beyond the current one-logarithm determinant
deficit.

\appendix

\section{Classical transcendence inputs and their exact roles}\label{app:classical}

This appendix states the external theorems used in the new paper proofs.  It is included to
make clear which conclusions are elementary consequences of the repository hypotheses and
which depend on major classical transcendence results.

\subsection{Gelfond--Schneider}

If $\alpha$ is algebraic with $\alpha\ne0,1$ and $z$ is algebraic irrational, then every value
of $\alpha^z$ is transcendental.  In this report it is used twice:

\begin{enumerate}
\item a nonintegral solution $x$ is first shown irrational and then must be transcendental
      because $2^x$ is algebraic;
\item $\vartheta=\log_2 3$ is irrational and must be transcendental because
      $2^{\vartheta}=3$.
\end{enumerate}

The theorem says nothing directly when the exponent is already transcendental, which is why
it cannot settle the original problem.

\subsection{Baker's homogeneous theorem}

A form sufficient here is: logarithms of algebraic numbers that are linearly independent over
$\Q$ are linearly independent over $\Qbar$.  It is used to upgrade the row independence in
\cref{thm:strong-spectrum} and the logarithmic independence needed in
\cref{thm:E3,thm:E2}.  The ordinary exact algebraic spectrum
\cref{thm:general-spectrum} does not require this upgrade; rational independence suffices for
six exponentials.

\subsection{Six exponentials}

If $u_1,u_2,u_3$ are linearly independent over $\Q$ and $v_1,v_2$ are linearly independent
over $\Q$, then at least one of
\[
  e^{u_iv_j}\qquad(1\le i\le3,\ 1\le j\le2)
\]
is transcendental.  Its uses are:

\begin{itemize}
\item exact algebraic-base spectrum: rows $\log a,\log b,\log\alpha$, columns $1,x$;
\item conditional algebraicity of $E_3$: rows $1,x,\vartheta$, columns $\log2,\log3$;
\item the symmetric $E_2$ argument;
\item the ordinary version of the unconditional $E_3/E_2$ dichotomy.
\end{itemize}

\subsection{Five exponentials}

For two independent pairs $u_1,u_2$ and $v_1,v_2$, and nonzero algebraic $\gamma$, at least
one of
\[
 e^{u_1v_1},e^{u_1v_2},e^{u_2v_1},e^{u_2v_2},
 e^{\gamma u_2/u_1}
\]
is transcendental.  It supplies the forbidden rays in \cref{prop:five-rays}.  It does not
remove the fifth exponential and hence does not prove four exponentials.

\subsection{Roy's strong six exponentials theorem}

Let $\LL$ be the $\Qbar$-linear span of $1$ and logarithms of algebraic numbers.  If
$u_1,u_2,u_3$ are independent over $\Qbar$ and $v_1,v_2$ are independent over $\Qbar$, then
one product $u_iv_j$ lies outside $\LL$.  This yields:

\begin{itemize}
\item the strong auxiliary-base spectrum;
\item the strong symmetric-constant dichotomy;
\item log-strong transcendence in the kernel and tower theorems.
\end{itemize}

Every ordinary transcendence conclusion in the report can be retained if this theorem is
removed; only the strengthened ``log outside $\LL$'' clauses disappear.

\subsection{Logical dependency graph}

\begin{center}
\begin{tabularx}{0.96\textwidth}{>{\bfseries\raggedright\arraybackslash}p{0.30\textwidth}>{\raggedright\arraybackslash}X}
\toprule
Result & Required input \\
\midrule
Nonintegral $x$ transcendental & Rational-integrality lemma + Gelfond--Schneider \\
Exact algebraic spectrum & Six exponentials \\
Strong spectrum & Exact row independence + Baker + Roy \\
Valuation lattice/Smith form & Exact spectrum + unique factorization \\
Kernel dimension and four types & Exact spectrum + elementary prime valuations \\
Kernel invariance & Exact spectrum applied twice \\
M\"obius criterion & Kernel definition + real logarithms \\
Depth-three theorem & Exact spectrum + kernel dimension; Roy only for strong form \\
Base-$2$ tower equivalence & Natural-base algebraicity classification + smooth branch algebra \\
Symmetric constants & Six exponentials; Baker/Roy for strong form \\
\bottomrule
\end{tabularx}
\end{center}

\section{Further consequences of the kernel formalism}\label{app:kernel-consequences}

\subsection{Orbit membership is global}

\begin{proposition}
Assume the conjecture fails.  Either every nonintegral solution belongs to
$\PGL_2(\Q)\cdot\vartheta$, or none does.
\end{proposition}

\begin{proof}
By \cref{thm:K-invariant}, all nonintegral solutions have the same kernel.  By
\cref{thm:mobius}, M\"obius-orbit membership is equivalent to nonzero kernel.
\end{proof}

This also follows from the exact solution monoid: if the least solution $\beta$ lies in the
orbit, then every $n+k\beta$ with $k>0$ is a rational affine transform of an orbit point and
therefore remains in the orbit.  Kernel invariance is stronger because it does not need the
least-generator theorem.

\subsection{The kernel as the unique rational relation space}

Define
\[
 R_x=\left\{(u,v,r,s)\in\Q^4:
  u+v\vartheta-rx-sx\vartheta=0
 \right\}.
\]
Projection to $(r,s)$ identifies $R_x$ with $K_x$: the coefficients $(u,v)$ are unique because
$1,\vartheta$ are independent.  Hence
\[
 \dim_\Q R_x=\dim_\Q K_x\in\{0,1\}.
\]
This gives a normalized linear-algebra version of the M\"obius criterion.  In Type IV the
four normalized logarithmic products
\[
  1,\quad \vartheta,\quad x,\quad x\vartheta
\]
are rationally independent even though the corresponding $2\times2$ real matrix has rank
one.  In Types I--III there is exactly one rational relation beyond the real determinant.

\subsection{Exact delayed-survival subgroup}

For a fixed nonintegral solution, define
\[
 \mathcal D_x=\{\alpha\in\Gamma:\alpha^{x^2}\in\Qbar\}.
\]
Then, under the coefficient identification $\Gamma\cong\Q^2$,
\[
  \mathcal D_x=K_x.
\]
Thus:

\begin{itemize}
\item in Type IV, every nontrivial base in $\Gamma$ fails by depth two;
\item in Types I--III, exactly one rational rank-one subgroup survives to depth two;
\item no nontrivial base survives to depth three while remaining algebraic at every prior
      level.
\end{itemize}

\subsection{Mixed second iterates as a quotient-rank invariant}

Let $\Apos/\Gamma$ be viewed as a $\Q$-vector space.  The two classes $[M],[A]$ span a
subspace of dimension one or two.  The kernel is the relation space of these classes:
\[
 K=\{(r,s):r[M]+s[A]=0\}.
\]
Therefore
\[
 \dim K=2-\dim\Span_\Q\{[M],[A]\}.
\]
Type III and the two axes have quotient rank one; Type IV has quotient rank two.  This
coordinate-free viewpoint explains why the kernel is solution-independent: the algebraicity
of $\alpha^{xy}$ detects the same quotient relation from either factor.

\section{Reproducibility inventory of the inspected corpus}\label{app:inventory}

\subsection{Article sources}

The three LaTeX blobs inspected at the repository commit on the title page were:

\begin{longtable}{@{}>{\raggedright\arraybackslash}p{0.59\textwidth}>{\raggedright\arraybackslash}p{0.33\textwidth}@{}}
\toprule
Path & Git blob SHA \\
\midrule
\endhead
\path{Article/alaoglu-erdos-combined-report/alaoglu-erdos-combined-report.tex}
& \sha{077723066e7a684a2c2cd1a7b0c49a02cfdf990b} \\
\path{Article/alaoglu-erdos-further-report/alaoglu-erdos-further-report.tex}
& \sha{d6bb12b447be426e675609b89a8b8b070e724b1e} \\
\path{Article/alaoglu-erdos-report-3/alaoglu-erdos-report-3.tex}
& \sha{048a3bcd902b5a09d06a68767e4aca7896aa0c33} \\
\bottomrule
\end{longtable}

The reports themselves cite earlier base commits for some theorem inventories.  This audit
instead compared their claims with the current aggregate import at
\sha{4848c6d834190301e014028a9c78a3f0ef49b7b8}.

\subsection{Current modules especially relevant to the new work}

The current branch includes the following modules, among others:

\begin{itemize}
\item \texttt{ArithmeticRigidity.lean};
\item \texttt{AlgebraicConsequences.lean};
\item \texttt{AlgebraicIntegerDeterminant.lean};
\item \texttt{AlgebraicOutputBridge.lean};
\item \texttt{AlgebraicThirdBase.lean};
\item \texttt{AlgebraicSixExponentials.lean};
\item \texttt{AlgebraicRationalBase.lean};
\item \texttt{AlgebraicBaseUnitPower.lean};
\item \texttt{AlgebraicRankThreeEquivalence.lean};
\item \texttt{AlgebraicOutputLocus.lean};
\item \texttt{CanonicalRadicalAlgebraicOutput.lean};
\item \texttt{CanonicalRadicalDivisibleHull.lean};
\item \texttt{CanonicalRadicalFieldNorm.lean};
\item \texttt{CanonicalRadicalSharpness.lean};
\item \texttt{CanonicalRadicalThirdOutput.lean};
\item \texttt{CanonicalRadicalValuation.lean};
\item \texttt{CoreIndependence.lean};
\item \texttt{DenominatorGrowth.lean};
\item \texttt{FiniteAlgebraicOutputBridge.lean};
\item \texttt{RationalBaseThird.lean};
\item \texttt{RationalSixExponentials.lean};
\item \texttt{RationalThirdBase.lean}.
\end{itemize}

In particular, the current snapshot now proves the exact $2,3$-divisible-hull locus for all
positive algebraic auxiliary bases, together with rank restrictions on mixed second
iterates.  This report generalizes the ordinary locus to an arbitrary independent algebraic
pair $(a,b)$, refines rational and integral outputs by a Smith lattice, and strengthens the
second-iterate slice to a solution-independent rational kernel and depth-three orbit theory.

\subsection{Representative theorem-status map}

\begin{longtable}{p{0.45\textwidth}p{0.16\textwidth}p{0.31\textwidth}}
\toprule
Statement & Status & Location/type \\
\midrule
\endhead
Nonintegral solution is transcendental & Lean & core transcendence modules \\
No candidate below $2^{12}$ & Lean & finite exclusion modules \\
Exact conditional solution monoid & Lean & conditional structure modules \\
Exact finite trigonometric block averages & Lean & block Weyl modules \\
Continuous-test-function equidistribution & Paper & combined report \\
Arbitrary-node Vandermonde repulsion & Paper & further report \\
$O((\log X)^2)$ dyadic candidate bound & Paper & further report \\
Nonlinear product criterion in $\Sol$ & Paper & third report \\
Primitive-solution quadratic asymptotic & Paper & third report \\
No candidate below $2^{26}$ & Computation & third-report MPFR source \\
Natural algebraic auxiliary-base classification & Lean & algebraic third-base modules \\
All positive algebraic auxiliary bases ($2,3$ locus) & Lean &
\lean{AlgebraicBaseUnitPower} \\
General arbitrary-pair saturated spectrum & New paper & \cref{thm:general-spectrum} \\
Strong $\LL$ spectrum & New paper & \cref{thm:strong-spectrum} \\
Valuation/Smith rational-output lattice & New paper & \cref{sec:valuation-spectrum} \\
Mixed squared-output rank-one slices & Lean & \lean{AlgebraicOutputLocus} \\
Global second-iterate kernel and four types & New paper & \cref{sec:kernel} \\
Depth-three tower theorem & New paper & \cref{sec:tower} \\
\bottomrule
\end{longtable}

\section{Common false shortcuts and audit warnings}\label{app:false-shortcuts}

The problem attracts several plausible but invalid one-line arguments.  Recording them helps
future formal and paper work avoid accidental circularity.

\begin{warning}[Fractional-part error]
Writing $x=n+\alpha$ with $0<\alpha<1$ and $2^x\in\Z$ does \emph{not} imply
$2^\alpha\in\Z$.  It implies only $2^\alpha\in\Z[1/2]$, and an irrational $\alpha$ is not an
ordinary radical exponent $1/q$.
\end{warning}

\begin{warning}[Rational-root circularity]
From $M^{\vartheta}\in\Z$ one cannot conclude that $\vartheta$ is rational or that $M$ is a
perfect power.  Gelfond--Schneider constrains algebraic irrational exponents, whereas
$\vartheta$ is transcendental.
\end{warning}

\begin{warning}[Baker is linear, not bilinear]
The relation
$\log M\log3=\log A\log2$ is not a Baker linear form.  Treating $\log2$ and $\log3$ as
algebraic coefficients is invalid.
\end{warning}

\begin{warning}[No modular real exponent]
The expression $2^x\bmod p$ has no canonical meaning for a general real $x$.  A $p$-adic
exponential requires a $p$-adic exponent and a unit in its convergence domain; neither is
supplied by the real equality.
\end{warning}

\begin{warning}[Candidate quotients need not be candidates]
If $m_1,m_2\in\Cand$ and $m_1\mid m_2$, it does not follow that $m_2/m_1\in\Cand$, because the
integer quotient of the outputs need not be integral.  Multiplication is safe; division
requires a separate valuation check.
\end{warning}

\begin{warning}[Density alone is compatible]
The orbit $\{kx\}$ is dense modulo one for irrational $x$, but the associated rational points
have exponential height.  An $O(\log H)$ sequence up to height $H$ is compatible with generic
transcendental-curve counting; density does not yield a contradiction without denominator
sensitivity.
\end{warning}

\begin{warning}[Numerical tolerance is not exclusion]
Computing $M^{\vartheta}$ in floating point and checking its distance to the nearest integer
against a fixed epsilon cannot certify a finite zero-free region.  Outward interval bounds or
an exact certificate are required.
\end{warning}

\begin{warning}[Algebraicity of a power is branch-sensitive]
From $z^d\in\Qbar$ with integer $d>0$ one may conclude $z\in\Qbar$, but from
$z^x\in\Qbar$ with transcendental $x$ no such algebraic closure argument is available.
\end{warning}

\section{Notation index}

\begin{longtable}{p{0.18\textwidth}p{0.72\textwidth}}
\toprule
Symbol & Meaning \\
\midrule
\endhead
$\Sol$ & Real $x$ with $2^x,3^x\in\Z$. \\
$\Cand$ & Positive integers $M$ with $M^{\vartheta}\in\N$. \\
$\vartheta$ & $\log3/\log2=\log_2 3$. \\
$\eta$ & $1/\vartheta=\log_3 2$. \\
$M,A$ & Integer outputs $M=2^x$, $A=3^x$. \\
$\beta$ & Conditional least nonintegral solution. \\
$\Apos$ & Positive real algebraic numbers. \\
$\GammaSat(a,b)$ & $\{a^rb^s:r,s\in\Q\}$. \\
$\Gamma$ & $\GammaSat(2,3)$. \\
$\LL$ & $\Qbar$-linear span of $1$ and logarithms of algebraic numbers. \\
$V_x$ & Full prime-valuation matrix with columns $\nu(M),\nu(A)$. \\
$\Lambda_x$ & Rational coefficient lattice giving rational auxiliary outputs. \\
$K_x$ & Kernel of external valuations, equivalently directions returning to $\Gamma$. \\
$E_3,E_2$ & $3^{\vartheta}$ and $2^{1/\vartheta}$. \\
\bottomrule
\end{longtable}

\begin{thebibliography}{99}

\bibitem{AlaogluErdos1944}
L. Alaoglu and P. Erd\H{o}s,
\emph{On highly composite and similar numbers},
Transactions of the American Mathematical Society \textbf{56} (1944), 448--469.
\url{https://doi.org/10.1090/S0002-9947-1944-0011087-2}.

\bibitem{Baker1975}
A. Baker,
\emph{Transcendental Number Theory},
Cambridge University Press, 1975.

\bibitem{Gelfond1960}
A. O. Gel'fond,
\emph{Transcendental and Algebraic Numbers},
Dover reprint, 1960.

\bibitem{KarlinStudden1966}
S. Karlin and W. J. Studden,
\emph{Tchebycheff Systems: With Applications in Analysis and Statistics},
Interscience, 1966.

\bibitem{Lang1966}
S. Lang,
\emph{Introduction to Transcendental Numbers},
Addison--Wesley, 1966.

\bibitem{Laurent2008}
M. Laurent,
\emph{Linear forms in two logarithms and interpolation determinants II},
Acta Arithmetica \textbf{133} (2008), no.~4, 325--348.

\bibitem{Ramachandra1968}
K. Ramachandra,
\emph{Contributions to the theory of transcendental numbers (I)},
Acta Arithmetica \textbf{14} (1968), no.~1, 65--72.

\bibitem{Roy1992}
D. Roy,
\emph{Matrices whose coefficients are linear forms in logarithms},
Journal of Number Theory \textbf{41} (1992), no.~1, 22--47.
\url{https://doi.org/10.1016/0022-314X(92)90081-Y}.

\bibitem{Waldschmidt2000}
M. Waldschmidt,
\emph{Diophantine Approximation on Linear Algebraic Groups: Transcendence Properties of the
Exponential Function in Several Variables},
Grundlehren der mathematischen Wissenschaften 326, Springer, 2000.

\bibitem{Waldschmidt2003}
M. Waldschmidt,
\emph{Open Diophantine Problems},
Moscow Mathematical Journal \textbf{4} (2004), no.~1, 245--305;
preprint \url{https://arxiv.org/abs/math/0312440}.

\bibitem{Waldschmidt2005}
M. Waldschmidt,
\emph{Variations on the Six Exponentials Theorem},
in R. Tandon (ed.), \emph{Algebra and Number Theory}, Hindustan Book Agency, 2005,
338--355.
\url{https://doi.org/10.1007/978-93-86279-23-1_23}.

\bibitem{Diaz2007}
G. Diaz,
\emph{Produits et quotients de combinaisons lin\'eaires de logarithmes de nombres
alg\'ebriques: conjectures et r\'esultats partiels},
Journal de Th\'eorie des Nombres de Bordeaux \textbf{19} (2007), no.~2, 373--391.
\url{https://doi.org/10.5802/jtnb.592}.

\bibitem{ProveItCombined}
The ProveIt project,
\emph{The Alaoglu--Erd\H{o}s Integer-Exponent Conjecture: Combined Report},
source blob \sha{077723066e7a684a2c2cd1a7b0c49a02cfdf990b}, 2026.

\bibitem{ProveItFurther}
The ProveIt project,
\emph{Further Progress on the Alaoglu--Erd\H{o}s Integer-Exponent Conjecture},
source blob \sha{d6bb12b447be426e675609b89a8b8b070e724b1e}, 2026.

\bibitem{ProveItReport3}
The ProveIt project,
\emph{Alaoglu--Erd\H{o}s Report 3},
source blob \sha{048a3bcd902b5a09d06a68767e4aca7896aa0c33}, 2026.

\bibitem{ProveItRepo}
V. Reshetnikov,
\emph{ProveIt}, GitHub repository, commit
\sha{4848c6d834190301e014028a9c78a3f0ef49b7b8}, 2026.
\url{https://github.com/VladimirReshetnikov/ProveIt}.

\bibitem{TheoremDBFourExp}
TheoremDB,
\emph{Four exponentials conjecture}, problem P4, reviewed July 31, 2026.
\url{https://theoremdb.org/statements/P4/}.

\end{thebibliography}

\end{document}
